% Photoelectric Effect — Physics Upper Sixth — Cours signé OnBuch+
% Official MINESEC syllabus. Compiler avec : tectonic PHY10-photoelectric-effect.tex
\newcommand{\DOCMATIERE}{Physics}
\newcommand{\DOCNIVEAU}{Upper Sixth}
\newcommand{\DOCMODULE}{Upper Sixth — Physics}
\newcommand{\DOCLECON}{10}
\newcommand{\DOCTITRE}{Photoelectric Effect}
% OnBuch+ — préambule « cours signés » ENRICHI (compilé avec tectonic / XeLaTeX)
% Système visuel « Pop » d'OnBuch+ : fond crème (jamais blanc), bordures encre
% épaisses, ombres « dures » décalées, titres Archivo Black en capitales.
% Version MATHÉMATIQUES : graphes (pgfplots/tikz), boîtes pédagogiques
% (définition, propriété, méthode, attention, le savais-tu, exercice résolu,
% exercice, TP, bilan), page de garde et signature OnBuch+ ancrée en bas.
%
% Macros attendues dans main.tex AVANT \input{preamble} :
%   \DOCMATIERE  \DOCNIVEAU  \DOCMODULE  \DOCLECON (numéro)  \DOCTITRE
\documentclass[11pt]{article}

\usepackage{fontspec}
\usepackage[english]{babel}
\usepackage[a4paper,margin=2cm,top=3.4cm,bottom=2.8cm,headheight=1.4cm,footskip=1.3cm]{geometry}
\usepackage{amsmath,amssymb,amsfonts}
\usepackage{mathtools} % \xmapsto, \coloneqq... souvent utilisés par les modèles
\usepackage{cancel} % simplifications barrées (bilans de forces)
% \abs{z} (module, valeur absolue) et \norm{u} (norme), utilisés par les modèles
\providecommand{\abs}{}\let\abs\relax\DeclarePairedDelimiter{\abs}{\lvert}{\rvert}
\providecommand{\norm}{}\let\norm\relax\DeclarePairedDelimiter{\norm}{\lVert}{\rVert}
\usepackage[table]{xcolor} % [table] : fournit aussi \rowcolors (alternance de lignes)
\usepackage{pagecolor}
\usepackage{tikz}
\usetikzlibrary{arrows.meta,positioning,calc,shapes.geometric,shapes.misc,shapes.arrows,decorations.pathmorphing,decorations.pathreplacing,decorations.markings,patterns,shadows,mindmap,trees,fit,backgrounds,matrix,intersections,angles,quotes}
\usepackage{pgfplots}
\usepackage[version=4]{mhchem} % équations nucléaires \ce{^{235}_{92}U}
\usepackage[european,siunitx]{circuitikz} % schémas électriques
\pgfplotsset{compat=1.18}
\usepgfplotslibrary{fillbetween,groupplots} % aires entre courbes (calcul intégral)
\usepackage{enumitem}
\usepackage{fancyhdr}
% [most] charge toutes les bibliothèques tcolorbox (listings, minted, raster,
% documentation...) et gonfle inutilement la mémoire TeX sur un document long
% (~300 boîtes) : on ne charge que ce qui est réellement utilisé.
\usepackage[skins,breakable]{tcolorbox}
\usepackage{graphicx}
\usepackage{colortbl}
\usepackage{booktabs}
\usepackage{tabularx}
\usepackage{diagbox} % case d'en-tête diagonale des tableaux à double entrée
% Colonnes extensibles centrée / gauche / droite (C, L, R), souvent utilisées par les modèles
\newcolumntype{C}{>{\centering\arraybackslash}X}
\newcolumntype{L}{>{\raggedright\arraybackslash}X}
\newcolumntype{R}{>{\raggedleft\arraybackslash}X}
\usepackage{array}
\usepackage{multicol}
\usepackage{siunitx}
% parse-numbers=false : accepte aussi \SI{2,5 \times 10^{-3}}{...}
\sisetup{locale=UK,output-decimal-marker={.},group-minimum-digits=4,parse-numbers=false}
\DeclareSIUnit\ohm{\ensuremath{\symup{\Omega}}} % U+2126 absent de la police
\DeclareSIUnit\division{div} % réglages d'oscilloscope (ms/div, V/div)
\DeclareSIUnit\tr{tr} % tours (vitesse de rotation, ex. \SI{1500}{\tr\per\minute})
\usepackage{qrcode}
% \degree : commande fréquemment utilisée par les modèles pour un angle ou une
% température en dehors de \SI{}{} (non fournie par les paquets chargés).
\providecommand{\degree}{^{\circ}}
% \qed (fin de démonstration), utilisé par les modèles sans amsthm
\providecommand{\qed}{\hfill$\square$}
% \barre : double barre des valeurs interdites dans un tableau de variations (tkz-tab)
\providecommand{\barre}{\ensuremath{\|}}
% environnement proof (amsthm non chargé) utilisé par les modèles
\ifdefined\proof\else\newenvironment{proof}[1][Proof]{\par\noindent\textit{#1.}\ }{\qed\par}\fi
\usepackage{lastpage}
\usepackage{hyperref}
\hypersetup{hidelinks}

% ============================================================
% Palette « Pop » OnBuch+ — mêmes hex que la classe P (plus_ui.dart)
% ============================================================
\definecolor{poporange}{HTML}{F59321}
\definecolor{popdark}{HTML}{E07A0C}
\definecolor{popcream}{HTML}{FFF6E8}
\definecolor{popink}{HTML}{1C1714}
\definecolor{poppurple}{HTML}{7A5AE0}
\definecolor{popgreen}{HTML}{2E9E6A}
\definecolor{poppink}{HTML}{E0457B}
\definecolor{popblue}{HTML}{2F7DE1}
\definecolor{popgold}{HTML}{C98A12}
\definecolor{popmuted}{HTML}{8A7F76}
\definecolor{popline}{HTML}{EADFD2}
\definecolor{popgray}{HTML}{8A7F76}
\colorlet{popgrey}{popgray}
% Alias tolérants (le modèle nomme parfois les couleurs différemment)
\colorlet{popwhite}{white}
\colorlet{popblack}{popink}
\colorlet{popred}{poppink}
\colorlet{popyellow}{popgold}
% Teintes claires pour fonds de tableaux / zones
\colorlet{popblueL}{popblue!10!white}
\colorlet{poporangeL}{poporange!14!white}
\colorlet{popgreenL}{popgreen!12!white}
\colorlet{poppurpleL}{poppurple!12!white}
\colorlet{poppinkL}{poppink!12!white}
\colorlet{popgoldL}{popgold!14!white}

\pagecolor{popcream}

% ============================================================
% Typographie
% ============================================================
\defaultfontfeatures{Ligatures=TeX}
\setmainfont{PlusJakartaSans-Medium.ttf}[Path=../../../fonts/,BoldFont=PlusJakartaSans-Bold.ttf]
% Maths en sans-serif (Fira Math) avec chiffres et lettres droites de Plus
% Jakarta, pour que les formules se fondent dans le texte.
\usepackage[math-style=ISO,bold-style=ISO]{unicode-math}
\setmathfont{FiraMath-Regular.otf}[Path=../../../fonts/]
\setmathfont{PlusJakartaSans-Medium.ttf}[Path=../../../fonts/,range={up/num,up/{Latin,latin},\mathup}]
% (icomma retiré : décimales avec point en anglais)
% Caractères mathématiques Unicode absents des polices de texte (Plus Jakarta,
% Archivo Black) : rendus par leur équivalent mathématique, en texte comme en
% formule — sinon ils s'affichent en carré vide (ex. « Arithmétique dans ℤ »).
\usepackage{newunicodechar}
\newunicodechar{ℝ}{\ensuremath{\mathbb{R}}}
\newunicodechar{ℤ}{\ensuremath{\mathbb{Z}}}
\newunicodechar{ℂ}{\ensuremath{\mathbb{C}}}
\newunicodechar{ℕ}{\ensuremath{\mathbb{N}}}
\newunicodechar{ℚ}{\ensuremath{\mathbb{Q}}}
\newunicodechar{𝔻}{\ensuremath{\mathbb{D}}}
\newunicodechar{α}{\ensuremath{\alpha}}
\newunicodechar{β}{\ensuremath{\beta}}
\newunicodechar{γ}{\ensuremath{\gamma}}
\newunicodechar{δ}{\ensuremath{\delta}}
\newunicodechar{ε}{\ensuremath{\varepsilon}}
\newunicodechar{θ}{\ensuremath{\theta}}
\newunicodechar{λ}{\ensuremath{\lambda}}
\newunicodechar{μ}{\ensuremath{\mu}}
\newunicodechar{σ}{\ensuremath{\sigma}}
\newunicodechar{τ}{\ensuremath{\tau}}
\newunicodechar{φ}{\ensuremath{\varphi}}
\newunicodechar{ψ}{\ensuremath{\psi}}
\newunicodechar{ω}{\ensuremath{\omega}}
\newunicodechar{Σ}{\ensuremath{\Sigma}}
% majuscules produites par \MakeUppercase dans les titres (α-aminés -> Α)
\newunicodechar{Α}{\ensuremath{\alpha}}
\newunicodechar{Β}{\ensuremath{\beta}}
\newunicodechar{Γ}{\ensuremath{\Gamma}}
\newunicodechar{Δ}{\ensuremath{\Delta}}
\newunicodechar{Ε}{\ensuremath{\varepsilon}}
\newunicodechar{Θ}{\ensuremath{\Theta}}
\newunicodechar{Λ}{\ensuremath{\Lambda}}
\newunicodechar{Μ}{\ensuremath{\mu}}
\newunicodechar{Τ}{\ensuremath{\tau}}
\newunicodechar{Φ}{\ensuremath{\Phi}}
\newunicodechar{Ψ}{\ensuremath{\Psi}}
\newunicodechar{Ω}{\ensuremath{\Omega}}
\newunicodechar{↦}{\ensuremath{\mapsto}}
\newunicodechar{∈}{\ensuremath{\in}}
\newunicodechar{∉}{\ensuremath{\notin}}
\newunicodechar{∧}{\ensuremath{\wedge}}
\newunicodechar{⇔}{\ensuremath{\Leftrightarrow}}
\newunicodechar{⟺}{\ensuremath{\Longleftrightarrow}}
\newunicodechar{⇒}{\ensuremath{\Rightarrow}}
\newunicodechar{⟹}{\ensuremath{\Longrightarrow}}
\newunicodechar{∘}{\ensuremath{\circ}}
\newunicodechar{∪}{\ensuremath{\cup}}
\newunicodechar{∩}{\ensuremath{\cap}}
\newunicodechar{≡}{\ensuremath{\equiv}}
\newunicodechar{⊂}{\ensuremath{\subset}}
\newunicodechar{⊥}{\ensuremath{\perp}}
\newunicodechar{∃}{\ensuremath{\exists}}
\newunicodechar{∀}{\ensuremath{\forall}}
\newunicodechar{∛}{\ensuremath{\sqrt[3]{\,}}}
\newunicodechar{∜}{\ensuremath{\sqrt[4]{\,}}}
\newunicodechar{⃗}{\ensuremath{{}^{\to}}}
\newunicodechar{ⁿ}{\ifmmode^{n}\else\textsuperscript{\ensuremath{n}}\fi}
\newunicodechar{ˣ}{\ifmmode^{x}\else\textsuperscript{\ensuremath{x}}\fi}
\newunicodechar{ᵇ}{\ifmmode^{b}\else\textsuperscript{\ensuremath{b}}\fi}
\newunicodechar{ʳ}{\ifmmode^{r}\else\textsuperscript{\ensuremath{r}}\fi}
\newunicodechar{ᵘ}{\ifmmode^{u}\else\textsuperscript{\ensuremath{u}}\fi}
\newunicodechar{ʸ}{\ifmmode^{y}\else\textsuperscript{\ensuremath{y}}\fi}
\newunicodechar{ᵃ}{\ifmmode^{a}\else\textsuperscript{\ensuremath{a}}\fi}
\newunicodechar{ᵉ}{\ifmmode^{e}\else\textsuperscript{\ensuremath{e}}\fi}
\newunicodechar{ᵏ}{\ifmmode^{k}\else\textsuperscript{\ensuremath{k}}\fi}
\newunicodechar{ᵖ}{\ifmmode^{p}\else\textsuperscript{\ensuremath{p}}\fi}
\newunicodechar{ᴺ}{\ifmmode^{N}\else\textsuperscript{\ensuremath{N}}\fi}
\newunicodechar{ᶜ}{\ifmmode^{c}\else\textsuperscript{\ensuremath{c}}\fi}
\newunicodechar{ʰ}{\ifmmode^{h}\else\textsuperscript{\ensuremath{h}}\fi}
\newunicodechar{ᵠ}{\ifmmode^{q}\else\textsuperscript{\ensuremath{q}}\fi}
\newunicodechar{ₙ}{\ifmmode_{n}\else\textsubscript{\ensuremath{n}}\fi}
\newunicodechar{ₐ}{\ifmmode_{a}\else\textsubscript{\ensuremath{a}}\fi}
\newunicodechar{ᵢ}{\ifmmode_{i}\else\textsubscript{\ensuremath{i}}\fi}
\newunicodechar{ᵣ}{\ifmmode_{r}\else\textsubscript{\ensuremath{r}}\fi}
\newunicodechar{ₖ}{\ifmmode_{k}\else\textsubscript{\ensuremath{k}}\fi}
\newunicodechar{ᵦ}{\ifmmode_{\beta}\else\textsubscript{\ensuremath{\beta}}\fi}
\newunicodechar{ₓ}{\ifmmode_{x}\else\textsubscript{\ensuremath{x}}\fi}
\newfontfamily\popdisplay{ArchivoBlack-Regular.ttf}[Path=../../../fonts/]
\newfontfamily\poplogo{SpaceGrotesk-Bold.ttf}[Path=../../../fonts/]
\newfontfamily\popsemi{PlusJakartaSans-SemiBold.ttf}[Path=../../../fonts/]
\newfontfamily\popxbold{PlusJakartaSans-ExtraBold.ttf}[Path=../../../fonts/]
\newfontfamily\popxboldls{PlusJakartaSans-ExtraBold.ttf}[Path=../../../fonts/,LetterSpace=3.0]

\setlist[enumerate]{leftmargin=1.7em,itemsep=5pt,topsep=4pt}
\setlist[itemize]{leftmargin=1.7em,itemsep=4pt,topsep=4pt,label={\color{poporange}\textbullet}}
\setlength{\parindent}{0pt}
\setlength{\parskip}{5pt}
\renewcommand{\arraystretch}{1.25}

% pgfplots : style Pop par défaut
\pgfplotsset{
  every axis/.append style={axis line style={popink,line width=1pt},tick style={popink},
    grid style={popline,line width=0.5pt},label style={font=\small},tick label style={font=\footnotesize}},
  popcurve/.style={poporange,line width=1.8pt,no markers,smooth},
}
% Même style disponible dans un \draw TikZ simple (hors pgfplots)
\tikzset{popcurve/.style={poporange,line width=1.8pt}}
\tikzset{popgrid/.style={popline,very thin}}
\colorlet{popcurve}{poporange} % le nom du style est parfois utilisé comme couleur

% ============================================================
% Pill / squiggle
% ============================================================
\newcommand{\poppill}[3]{%
  \tikz[baseline=-0.55ex]\node[draw=popink,line width=1.2pt,rounded corners=2.6pt,%
    fill=#2,inner xsep=6.5pt,inner ysep=3pt]{%
    \popxboldls\fontsize{7}{7}\selectfont\color{#3}\MakeUppercase{#1}};%
}
\newcommand{\popsquiggle}[1]{%
  \tikz[baseline=-0.4ex]{
    \draw[#1,line width=2.6pt,line cap=round]
      plot [smooth] coordinates {(0,0.09) (0.34,0.01) (0.68,0.21) (1.02,0.01) (1.36,0.10)};
  }%
}
% Mot-clé surligné (marqueur orange sous le texte)
\newcommand{\cle}[1]{{\popxbold\color{popdark}#1}}

% ============================================================
% En-tête / pied
% ============================================================
\pagestyle{fancy}
\fancyhf{}
\renewcommand{\headrulewidth}{0pt}
\renewcommand{\footrulewidth}{0pt}
\lhead{\raisebox{-0.15ex}{\poplogo\fontsize{13}{13}\selectfont\color{popink}OnBuch\color{poporange}+}}
\rhead{\poppill{\DOCMATIERE\ \textbullet\ \DOCNIVEAU\ \textbullet\ Chapter \DOCLECON}{white}{popink}}
\lfoot{\popxbold\fontsize{7.6}{7.6}\selectfont\color{popmuted}\MakeUppercase{Signed course OnBuch+}\ \textbullet\ {\popsemi\DOCTITRE}}
\rfoot{\tikz[baseline=-0.6ex]\node[rounded corners=8.5pt,fill=popink,minimum height=17pt,minimum width=17pt,inner xsep=5pt,inner sep=0]{\popxbold\fontsize{8}{8}\selectfont\color{popcream}\thepage\,/\,\pageref*{LastPage}};}

% ============================================================
% Titres
% ============================================================
\newcommand{\chaptitre}[2]{%
  \begin{tcolorbox}[enhanced,colback=poporange,colframe=popink,boxrule=2.6pt,arc=11pt,
      drop shadow=popink,left=15pt,right=15pt,top=12pt,bottom=11pt,before skip=3pt,after skip=18pt]
    \poppill{#2}{popink}{popcream}\par\vspace{9pt}
    {\raggedright\hyphenpenalty=10000\exhyphenpenalty=10000\popdisplay\fontsize{22}{24}\selectfont\color{popcream}\MakeUppercase{#1}\par}
  \end{tcolorbox}
}
% Section numérotée : pastille orange avec le numéro + titre Archivo
\newcounter{popsec}
\newcommand{\coursec}[1]{%
  \stepcounter{popsec}\setcounter{popsub}{0}%
  \par\vspace{16pt}\noindent
  \begin{minipage}{\linewidth}
  \tikz[baseline=-0.7ex]\node[draw=popink,line width=1.6pt,fill=poporange,rounded corners=4pt,minimum size=22pt,inner sep=0]{\popdisplay\fontsize{12}{12}\selectfont\color{popcream}\thepopsec};%
  \hspace{8pt}{\popdisplay\fontsize{15}{17}\selectfont\color{popink}\MakeUppercase{#1}}\par
  \vspace{4pt}\noindent{\color{poporange}\rule{36pt}{3.6pt}}
  \end{minipage}\par\nopagebreak\vspace{8pt}
}
\newcounter{popsub}
\newcommand{\courssub}[1]{%
  \stepcounter{popsub}%
  \par\vspace{9pt}\noindent{\popxbold\fontsize{11.5}{13}\selectfont\color{popink}{\color{poporange}\thepopsec.\thepopsub}\hspace{6pt}#1}\par\nopagebreak\vspace{4pt}
}

% ============================================================
% Boîtes pédagogiques — même recette : fond blanc, bordure épaisse colorée,
% ombre dure encre, badge plein. Titre optionnel [#1].
% ============================================================
\tcbset{popbase/.style={breakable,enhanced,colback=white,boxrule=2.3pt,arc=9pt,
  shadow={3pt}{-3pt}{0pt}{popink},left=14pt,right=14pt,top=11pt,bottom=11pt,before skip=11pt,after skip=15pt,
  fontupper=\popsemi\fontsize{10.6}{14.5}\selectfont\color{popink},
  fontlower=\popsemi\fontsize{10.6}{14.5}\selectfont\color{popink}}}
% Coche dessinée (le glyphe ✓ n'existe pas dans les polices du document)
\newcommand{\popcheck}[1]{\tikz[baseline=-0.5ex]\draw[#1,line width=1.6pt,line cap=round,line join=round](-0.13,0.0)--(-0.04,-0.09)--(0.13,0.1);}
\providecommand{\coche}{\popcheck{popgreen}}
\providecommand{\resultat}[1]{\par\smallskip\noindent\fcolorbox{popgreen}{popgreenL}{\parbox{\dimexpr\linewidth-2\fboxsep-2\fboxrule\relax}{\textbf{Result.} #1}}\par\smallskip}
\newcommand{\popboxhead}[3]{\poppill{#1}{#2}{white}\ifx\relax#3\relax\else\hspace{6pt}{\popxbold\fontsize{10.6}{12}\selectfont\color{popink}#3}\fi\par\vspace{7pt}}

\newtcolorbox{aretenir}[1][]{popbase,colframe=popgreen,before upper={\popboxhead{Key points}{popgreen}{#1}}}
\newtcolorbox{exemplebox}[1][]{popbase,colframe=poporange,before upper={\popboxhead{Exemple}{poporange}{#1}}}
\newtcolorbox{definition}[1][]{popbase,colframe=popblue,before upper={\popboxhead{Definition}{popblue}{#1}}}
\newtcolorbox{propriete}[1][]{popbase,colframe=poppurple,before upper={\popboxhead{Property}{poppurple}{#1}}}
\newtcolorbox{methode}[1][]{popbase,colframe=popgold,before upper={\popboxhead{Method}{popgold}{#1}}}
\newtcolorbox{attention}[1][]{popbase,colframe=poppink,before upper={\popboxhead{Warning}{poppink}{#1}}}
\newtcolorbox{experience}[1][]{popbase,colframe=popblue,colback=popblueL,before upper={\popboxhead{Experiment}{popblue}{#1}}}
% « Le savais-tu ? » : carte violette pleine (contexte, culture, Cameroun)
\newtcolorbox{savaistu}[1][]{popbase,colback=poppurple,colframe=popink,
  fontupper=\popsemi\fontsize{10.4}{14.2}\selectfont\color{white},
  before upper={\poppill{Did you know?}{white}{poppurple}\ifx\relax#1\relax\else\hspace{6pt}{\popxbold\color{white}#1}\fi\par\vspace{7pt}}}
% Exercice résolu : énoncé en haut, \tcblower puis la solution sur fond crème
\newcounter{popexo}\newcounter{popexr}
\newtcolorbox{exoresolu}[1][]{popbase,colframe=poporange,colbacklower=popcream,
  segmentation style={popink,line width=1.2pt,dashed},
  before upper={\stepcounter{popexr}\popboxhead{Worked example \thepopexr}{poporange}{#1}},
  before lower={\poppill{Solution}{popink}{popcream}\par\vspace{6pt}}}
% Exercice (non résolu en place) — la correction est dans la partie Corrigés
\newtcolorbox{exercice}[1][]{popbase,colframe=popink,
  before upper={\stepcounter{popexo}\popboxhead{Exercise \thepopexo}{popink}{#1}}}
% Corrigé
\newtcolorbox{corrige}[1][]{popbase,colframe=popgreen,colback=popgreenL,
  before upper={\popboxhead{Solution}{popgreen}{#1}}}
% Objectifs / prérequis / bilan
\newtcolorbox{objectifs}{popbase,colframe=popink,colback=poporangeL,
  before upper={\popboxhead{Chapter objectives}{popink}{}}}
\newtcolorbox{prerequis}{popbase,colframe=popink,colback=popblueL,
  before upper={\popboxhead{Prerequisites}{popblue}{}}}
\newtcolorbox{bilan}{popbase,colframe=popink,colback=white,boxrule=2.8pt,
  before upper={\popboxhead{Fiche bilan}{poporange}{L'essentiel à connaître pour l'examen}}}
% Figure encadrée avec légende
\newtcolorbox{popfigure}[1][]{enhanced,breakable=false,colback=white,colframe=popline,boxrule=1.4pt,arc=8pt,
  left=10pt,right=10pt,top=10pt,bottom=8pt,before skip=10pt,after skip=12pt,halign=center,
  lower separated=false,halign lower=center,
  fontlower=\popsemi\fontsize{9}{11.5}\selectfont\color{popmuted},
  #1}
\newcommand{\legende}[1]{\par\vspace{6pt}{\popsemi\fontsize{9}{11.5}\selectfont\color{popmuted}#1}}

% ============================================================
% Signature OnBuch+
% ============================================================
\newcommand{\poplogoinline}[1]{{\poplogo\fontsize{#1}{#1}\selectfont\color{popink}OnBuch\color{poporange}+}}
\newcommand{\signatureonbuch}{%
  \begin{tcolorbox}[enhanced,colback=popink,colframe=popink,boxrule=2.2pt,arc=10pt,
      drop shadow=poporange,left=14pt,right=14pt,top=10pt,bottom=10pt,before skip=0pt,after skip=0pt]
    \tikz[baseline=-0.7ex]\node[circle,fill=popgreen,draw=popcream,line width=1.3pt,minimum size=22pt,inner sep=0]{\popcheck{white}};%
    \hspace{9pt}\begin{minipage}[c]{0.62\linewidth}
      {\popxbold\fontsize{12}{13}\selectfont\color{popcream}Signed\ \textbullet\ {\poplogo OnBuch\color{poporange}+}}\par\vspace{2pt}
      {\raggedright\popsemi\fontsize{8}{10.5}\selectfont\color{popline}\MakeUppercase{OnBuch+ teaching team}\\\MakeUppercase{Follows the official MINESEC syllabus}\par}
    \end{minipage}\hfill
    {\poplogo\fontsize{22}{22}\selectfont\color{popcream}OnBuch\color{poporange}+}
  \end{tcolorbox}%
}

% ============================================================
% Page de garde
% #1 titre · #2 matière · #3 niveau · #4 module · #5 n° leçon · #6 durée ·
% #7 description / objectifs (texte ou itemize)
% ============================================================
\newcommand{\pagegarde}[7]{%
  \thispagestyle{empty}
  \newgeometry{a4paper,margin=1.7cm,top=1.6cm,bottom=1.4cm}%
  \begin{tikzpicture}[remember picture,overlay]
    \fill[poporange!18] ([xshift=-3.2cm,yshift=-3.6cm]current page.north east) circle (4.6cm);
    \fill[poppurple!14] ([xshift=2.4cm,yshift=7.6cm]current page.south west) circle (3.8cm);
  \end{tikzpicture}%
  {\poplogo\fontsize{30}{30}\selectfont\color{popink}OnBuch\color{poporange}+}\hfill
  \raisebox{0.6ex}{\poppill{Signed course \textbullet\ Premium}{popink}{popcream}}\par
  \vspace{1pt}\popsquiggle{poppurple}\par
  \vspace{8pt}
  \begin{tcolorbox}[enhanced,colback=poporange,colframe=popink,boxrule=2.8pt,arc=13pt,
      drop shadow=popink,left=18pt,right=18pt,top=12pt,bottom=13pt,after skip=10pt]
    \poppill{#2 \textbullet\ #3}{popink}{popcream}\hspace{5pt}\poppill{#4}{white}{popink}\par\vspace{8pt}
    {\popdisplay\fontsize{14}{14}\selectfont\color{popink}CHAPTER #5}\par\vspace{4pt}
    {\raggedright\hyphenpenalty=10000\exhyphenpenalty=10000\popdisplay\fontsize{25}{27}\selectfont\color{popcream}\MakeUppercase{#1}\par}
    \vspace{8pt}
    {\popxbold\fontsize{9.5}{9.5}\selectfont\color{popink}\MakeUppercase{Suggested time: #6}}
  \end{tcolorbox}
  \begin{tcolorbox}[enhanced,colback=white,colframe=popink,boxrule=2.2pt,arc=10pt,
      drop shadow=popink,left=16pt,right=16pt,top=10pt,bottom=10pt,after skip=8pt]
    \poppill{In this chapter}{poppurple}{white}\par\vspace{5pt}
    \popsemi\fontsize{9.3}{12.6}\selectfont\color{popink}#7
  \end{tcolorbox}
  \noindent\begin{minipage}[t]{0.487\linewidth}
    \begin{tcolorbox}[enhanced,colback=popgreen,colframe=popink,boxrule=2.2pt,arc=10pt,
        drop shadow=popink,left=11pt,right=11pt,top=10pt,bottom=10pt,height=3.1cm,valign=center]
      \tcbox[colback=white,colframe=popink,boxrule=1.3pt,arc=5pt,boxsep=0pt,nobeforeafter,
        left=3pt,right=3pt,top=3pt,bottom=3pt,valign=center]{\qrcode[height=1.9cm]{https://play.google.com/store/apps/details?id=com.luvvix.onbuch&hl=en}}%
      \hspace{8pt}\begin{minipage}[c]{0.5\linewidth}
        {\popdisplay\fontsize{11}{12.5}\selectfont\color{white}\MakeUppercase{Download OnBuch}}\par\vspace{4pt}
        {\raggedright\popsemi\fontsize{8.4}{11}\selectfont\color{white}Free \textbullet\ Google Play\\AI tutor, past papers, results\par}
      \end{minipage}
    \end{tcolorbox}
  \end{minipage}\hfill
  \begin{minipage}[t]{0.487\linewidth}
    \begin{tcolorbox}[enhanced,colback=poppurple,colframe=popink,boxrule=2.2pt,arc=10pt,
        drop shadow=popink,left=11pt,right=11pt,top=10pt,bottom=10pt,height=3.1cm,valign=center]
      \tikz[baseline=-0.7ex]\node[circle,fill=white,draw=popink,line width=1.3pt,minimum size=1.6cm,inner sep=0]{\popdisplay\fontsize{24}{24}\selectfont\color{poppurple}+};%
      \hspace{8pt}\begin{minipage}[c]{0.6\linewidth}
        {\popdisplay\fontsize{11}{12.5}\selectfont\color{white}\MakeUppercase{Go OnBuch+}}\par\vspace{4pt}
        {\raggedright\popsemi\fontsize{8.4}{11}\selectfont\color{white}All signed courses, unlimited AI tutor, PDF sheets — OnBuch+ tab in the app\par}
      \end{minipage}
    \end{tcolorbox}
  \end{minipage}
  \vspace{8pt}
  \popcontenu
  \vfill
  \signatureonbuch
  \restoregeometry
  \newpage
}

% Rangée « Ce cours contient » (page de garde)
\newcommand{\poptile}[3]{% #1 couleur #2 gros texte #3 légende
  \begin{tcolorbox}[enhanced,colback=#1,colframe=popink,boxrule=2pt,arc=9pt,shadow={2.5pt}{-2.5pt}{0pt}{popink},
      left=6pt,right=6pt,top=9pt,bottom=9pt,halign=center,valign=center,height=1.75cm,width=\linewidth,nobeforeafter]
    {\popdisplay\fontsize{12.5}{13}\selectfont\color{white}\MakeUppercase{#2}}\par\vspace{3pt}
    {\centering\popsemi\fontsize{7.8}{9.5}\selectfont\color{white}#3\par}
  \end{tcolorbox}}
\newcommand{\popcontenu}{%
  \par\noindent{\popxboldls\fontsize{7.5}{7.5}\selectfont\color{popmuted}THIS SIGNED COURSE CONTAINS}\par\vspace{7pt}
  \noindent\begin{minipage}[t]{0.235\linewidth}\poptile{popblue}{Course}{complete, illustrated, step by step}\end{minipage}\hfill
  \begin{minipage}[t]{0.235\linewidth}\poptile{poporange}{Methods}{and worked examples}\end{minipage}\hfill
  \begin{minipage}[t]{0.235\linewidth}\poptile{poppink}{Exercises}{exam style + solutions}\end{minipage}\hfill
  \begin{minipage}[t]{0.235\linewidth}\poptile{popgreen}{Summary sheet}{the essentials to remember}\end{minipage}\par}

% Sommaire
\newtcolorbox{sommaire}{enhanced,colback=white,colframe=popink,boxrule=2.2pt,arc=10pt,
  drop shadow=popink,left=18pt,right=18pt,top=13pt,bottom=13pt,before skip=6pt,after skip=20pt,
  before upper={\poppill{Contents}{poppurple}{white}\par\vspace{9pt}\popsemi\fontsize{10.6}{14.5}\selectfont\color{popink}}}

% Signature de fin de cours — poussée tout en bas de la dernière page
\newcommand{\findecours}{%
  \par\vfill\nopagebreak
  \begin{center}\popsquiggle{poporange}\end{center}\vspace{4pt}
  \signatureonbuch
}
\AtBeginDocument{\def\llbracket{\mathopen{[\mkern-3.5mu[}}\def\rrbracket{\mathclose{]\mkern-3.5mu]}}} % intervalles d'entiers (glyphe ⟦ absent de la police)
% Glyphes absents de Fira Math : redéfinis à partir de glyphes disponibles
\AtBeginDocument{%
  \def\setminus{\mathbin{\backslash}}\let\smallsetminus\setminus
  \def\vdots{\mathord{\vbox{\baselineskip3.2pt\lineskiplimit0pt\kern4pt\hbox{.}\hbox{.}\hbox{.}}}}%
  \def\ddots{\mathinner{\mkern1mu\raise6.5pt\hbox{.}\mkern2mu\raise3.6pt\hbox{.}\mkern2mu\raise0.7pt\hbox{.}\mkern1mu}}%
  \def\triangle{\mathord{\scalebox{0.9}{$\symup{\Delta}$}}}%
  \def\nabla{\mathord{\raisebox{\depth}{\scalebox{1}[-1]{$\symup{\Delta}$}}}}%
  \def\checkmark{\popcheck{popdark}}%
  \def\star{\mathbin{\ast}}\def\bigstar{\mathord{\ast}}%
  \def\diamond{\mathbin{\scalebox{0.6}{$\Diamond$}}}%
  \def\bigcup{\mathop{\mathchoice{\raisebox{-0.3em}{\scalebox{1.7}{$\cup$}}}{\scalebox{1.25}{$\cup$}}{\cup}{\cup}}\displaylimits}%
  \def\bigcap{\mathop{\mathchoice{\raisebox{-0.3em}{\scalebox{1.7}{$\cap$}}}{\scalebox{1.25}{$\cap$}}{\cap}{\cap}}\displaylimits}%
  \def\lesseqqgtr{\mathrel{\lessgtr}}\def\bigoplus{\mathop{\oplus}\displaylimits}%
}
\providecommand{\ssi}{if and only if} % abréviation fréquente
\AtBeginDocument{\providecommand{\wideparen}{\overparen}} % arc de cercle

\begin{document}

\pagegarde{Photoelectric Effect}{Physics}{Upper Sixth}{Upper Sixth — Physics}{10}{2 periods}{This chapter introduces the photoelectric effect, the photon model of light and Planck's constant, leading to Einstein's photoelectric equation. It ends with the wave-particle duality of light, a cornerstone of modern physics examined at the GCE Advanced Level.
\vspace{4pt}
\begin{multicols}{2}\small
\begin{itemize}
\item By the end of this chapter I can describe the photoelectric effect and state the conditions under which it occurs.
\item By the end of this chapter I can explain why the classical wave theory of light fails to account for the experimental observations.
\item By the end of this chapter I can define a photon and calculate its energy using E = hf and E = hc/λ.
\item By the end of this chapter I can use Planck's constant h = 6.63 × 10⁻³⁴ J·s correctly in calculations, including conversions between joules and electron-volts.
\item By the end of this chapter I can state and apply Einstein's photoelectric equation hf = W₀ + Ek(max).
\item By the end of this chapter I can determine the work function, threshold frequency and threshold wavelength of a metal from data or from a graph.
\end{itemize}\end{multicols}}

\begin{sommaire}
\begin{multicols}{2}
\begin{enumerate}
\item The Photoelectric Effect: Observation and Failure of Wave Theory
\item Photons and Planck's Constant
\item Einstein's Photoelectric Equation
\item Wave-Particle Duality and Applications
\item Integration activity
\item Methods and worked examples
\item Exercises
\item Solutions to the exercises
\item Summary sheet
\end{enumerate}
\end{multicols}
\end{sommaire}

\begin{objectifs}
\begin{itemize}
\item By the end of this chapter I can describe the photoelectric effect and state the conditions under which it occurs.
\item By the end of this chapter I can explain why the classical wave theory of light fails to account for the experimental observations.
\item By the end of this chapter I can define a photon and calculate its energy using E = hf and E = hc/λ.
\item By the end of this chapter I can use Planck's constant h = 6.63 × 10⁻³⁴ J·s correctly in calculations, including conversions between joules and electron-volts.
\item By the end of this chapter I can state and apply Einstein's photoelectric equation hf = W₀ + Ek(max).
\item By the end of this chapter I can determine the work function, threshold frequency and threshold wavelength of a metal from data or from a graph.
\item By the end of this chapter I can interpret the graph of maximum kinetic energy of photoelectrons against frequency of incident radiation.
\item By the end of this chapter I can explain the meaning of wave-particle duality and give evidence for both aspects of light.
\item By the end of this chapter I can relate the photoelectric effect to applications such as photocells, solar panels and light sensors.
\end{itemize}
\end{objectifs}

\begin{prerequis}
\begin{itemize}
\item Wave properties of light: frequency, wavelength and the relation c = fλ (Lower Sixth)
\item The electromagnetic spectrum and orders of magnitude of wavelengths
\item Energy, work and the electron-volt as a unit of energy
\item Structure of the atom: electrons, ionisation and energy levels (Lower Sixth)
\item Basic graph skills: gradient, intercepts and equation of a straight line
\end{itemize}
\end{prerequis}

\newpage

\coursec{The Photoelectric Effect: Observation and Failure of Wave Theory}

In Douala, solar panels on rooftops convert sunlight directly into electricity, and automatic street lights in Yaoundé switch on by themselves at dusk. Yet for decades, physicists could not explain why shining a faint blue light on a metal ejects electrons instantly, while an intense red light --- however bright --- ejects nothing at all. This puzzle, which classical wave theory completely failed to solve, forced Einstein to propose that light itself arrives in tiny packets of energy called photons. In this chapter we uncover how this revolutionary idea explains the photoelectric effect and opened the door to quantum physics.

\textbf{The problem of this section.} We begin with the experiment itself: what exactly happens when light falls on a metal? We describe the apparatus, state the four experimental laws that any correct theory must explain, and then show, point by point, why the classical wave picture of light --- so successful for interference and diffraction --- collapses completely here. Finally we introduce the \cle{stopping potential}, the experimental tool that lets us measure the energy of the ejected electrons.

\courssub{Observation: Light Can Tear Electrons Out of a Metal}

\textbf{Intuition first.} Inside a metal, some electrons (the conduction electrons) move almost freely between the atoms, but they are still bound to the metal as a whole: to escape through the surface, an electron must receive a minimum amount of energy. Heating the metal can supply this energy (thermionic emission, used in old radio valves). The surprising discovery, made by Hertz in 1887 while experimenting with electromagnetic waves, and studied in detail by Hallwachs and Lenard, is that \emph{light alone}, at room temperature, can also do the job --- but only under very strange conditions.

\begin{definition}[Photoelectric effect and photoelectrons]
The \cle{photoelectric effect} is the emission of electrons from the surface of a metal (or, more generally, a material) when electromagnetic radiation of sufficiently high frequency falls on it.

The emitted electrons are called \cle{photoelectrons}. They are perfectly ordinary electrons; the prefix ``photo'' simply recalls that they were set free by light.
\end{definition}

\begin{attention}[Photoelectrons are not a new particle]
A photoelectron has exactly the same charge $-e = -1.60\times 10^{-19}\ \mathrm{C}$ and the same mass $m_e = 9.11\times 10^{-31}\ \mathrm{kg}$ as any other electron. Only its \emph{origin} (ejection by light) is special.
\end{attention}

\courssub{The Experimental Set-Up: The Photoelectric Cell}

To study the effect quantitatively we use a \cle{photoelectric cell} (photoemissive cell). Its essential parts are:

\begin{itemize}
\item an \textbf{evacuated glass tube}, so that the ejected electrons travel without colliding with air molecules;
\item a metal plate, the \textbf{cathode} (or photocathode), coated with the metal under study (e.g.\ zinc, sodium, caesium), on which the light falls;
\item a collecting electrode, the \textbf{anode}, facing the cathode;
\item an external circuit containing a \textbf{variable potential difference} (potentiometer arrangement, allowing the anode to be made positive \emph{or} negative relative to the cathode), a \textbf{microammeter} (or sensitive galvanometer) to detect the tiny photoelectric current, and a voltmeter.
\end{itemize}

\begin{popfigure}
\begin{center}
\begin{tikzpicture}[scale=0.95, every node/.style={font=\small}]
% Glass envelope
\draw[thick, popblue] (0,0) circle (2.6);
\node[popblue] at (0,2.95) {evacuated tube};
% Cathode
\draw[very thick, popdark] (-1.6,-1.2) -- (-1.6,1.2);
\draw[very thick, popdark] (-1.6,0) -- (-3.4,0);
\node[align=center] at (-2.6,-1.7) {cathode C\\ (metal)};
% Anode
\draw[very thick, poporange] (1.3,-0.9) -- (1.3,0.9);
\draw[very thick, poporange] (1.3,0) -- (3.4,0);
\node[align=center, poporange] at (2.3,-1.4) {anode A};
% Light arrows
\foreach \y in {-0.7,0,0.7}{
  \draw[->, very thick, popgold] (-5.6,\y) -- (-1.9,\y);
}
\node[popgold] at (-4.4,1.15) {incident light};
% Electron paths
\foreach \y in {-0.5,0.1,0.6}{
  \draw[->, dashed, poppurple] (-1.45,\y) .. controls (0,\y+0.35) .. (1.15,\y+0.1);
}
\node[poppurple] at (0,1.35) {$e^-$};
% External circuit
\draw[thick] (-3.4,0) -- (-3.4,-3.4) -- (-1.2,-3.4);
\draw[thick] (3.4,0) -- (3.4,-3.4) -- (1.2,-3.4);
% Microammeter
\draw[thick] (-1.2,-3.4) circle (0.45);
\node at (-1.2,-3.4) {$\mu$A};
% Voltmeter
\draw[thick] (1.2,-3.4) circle (0.45);
\node at (1.2,-3.4) {V};
% Variable supply
\draw[thick] (-0.75,-3.4) -- (-0.45,-3.4);
\draw[thick] (-0.6,-3.4) -- (-0.6,-3.15);
\draw[thick] (-0.6,-3.65) -- (-0.6,-3.4);
\draw[thick] (0.75,-3.4) -- (0.45,-3.4);
\draw[thick] (0.6,-3.4) -- (0.6,-3.15);
\draw[thick] (0.6,-3.65) -- (0.6,-3.4);
\draw[thick] (-0.45,-3.4) -- (0.45,-3.4);
\draw[->, thick, popgreen] (-0.15,-4.05) -- (0.35,-3.45);
\node[align=center] at (0,-4.5) {variable p.d.\ $U_{AC}$ (reversible)};
\end{tikzpicture}
\end{center}
\legende{The photoelectric cell. Light falling on the cathode releases photoelectrons; those reaching the anode constitute the photoelectric current measured by the microammeter.}
\end{popfigure}

\textbf{How the measurement works.} When the anode is made \emph{positive} with respect to the cathode ($U_{AC} > 0$), the ejected electrons are attracted to the anode and a current $I$ flows, detected by the microammeter. The photoelectric current is simply the rate of flow of photoelectrons:
\[
I = \frac{Q}{t} = \frac{N e}{t},
\]
where $N$ is the number of photoelectrons collected in time $t$. Hence the current directly counts the emitted electrons: $N/t = I/e$. When $U_{AC}$ is made large enough, \emph{every} emitted electron is collected: the current then reaches a maximum value called the \cle{saturation current} $I_s$, which measures the total rate of emission from the cathode.

\courssub{The Four Experimental Laws of Photoelectric Emission}

Careful experiments (notably by Lenard, and later the precision measurements of Millikan) established four facts that any theory of the effect must explain.

\begin{propriete}[Experimental laws of the photoelectric effect]
\begin{enumerate}
\item \textbf{Threshold frequency.} For each metal there exists a minimum frequency $f_0$, called the \cle{threshold frequency}, below which \emph{no} photoelectron is emitted, whatever the intensity of the light. The value of $f_0$ is characteristic of the metal (equivalently, a threshold wavelength $\lambda_0 = c/f_0$).
\item \textbf{Instantaneous emission.} When $f > f_0$, emission begins essentially instantaneously (delay less than $10^{-9}\ \mathrm{s}$), even for extremely weak light.
\item \textbf{Current proportional to intensity.} For $f > f_0$, the number of photoelectrons emitted per second --- and therefore the saturation current --- is directly proportional to the intensity of the radiation.
\item \textbf{Kinetic energy depends on frequency only.} The maximum kinetic energy $E_{k(\max)}$ of the photoelectrons increases with the \emph{frequency} of the radiation and is completely \emph{independent of its intensity}.
\end{enumerate}
\end{propriete}

\begin{exemplebox}[Threshold frequencies of some metals]
Zinc emits electrons only under ultraviolet light ($f_0 \approx 8.1\times 10^{14}\ \mathrm{Hz}$, $\lambda_0 \approx 370\ \mathrm{nm}$), which is why a charged zinc plate under ordinary visible light keeps its charge. Sodium and caesium have lower thresholds and respond to visible light: caesium can even emit under red light. This metal-dependence of $f_0$ is itself a crucial clue: the threshold is a property of the \emph{material}, not of the light alone.
\end{exemplebox}

\courssub{What Classical Wave Theory Predicts --- and Why It Fails}

\textbf{The classical picture.} In the wave theory of the 19th century, light is a continuous electromagnetic wave. Its energy is spread smoothly over the whole wavefront, and the energy delivered per second to a given area is measured by the \emph{intensity}. An electron at the metal surface should therefore absorb energy gradually, like a small boat rocked by ocean swells, until it has accumulated enough to escape.

From this picture, three clear predictions follow:

\begin{enumerate}
\item \textbf{Any frequency should work.} Given enough time, a wave of any frequency delivers arbitrarily large energy. A feeble red light, shining long enough, should eventually eject electrons. There should be \emph{no threshold frequency}.
\item \textbf{A time delay should exist for weak light.} With very dim light the energy trickles in slowly, so electrons should be emitted only after a measurable delay --- seconds, minutes or even hours for the faintest sources.
\item \textbf{Kinetic energy should grow with intensity.} A more intense wave carries more energy per second, so it should shake the electrons harder and eject them with greater kinetic energy, whatever the frequency.
\end{enumerate}

\begin{center}
\begin{tabularx}{\linewidth}{|l|X|X|}
\hline
\rowcolor{popblueL} \textbf{Observation} & \textbf{Wave theory prediction} & \textbf{Experiment} \\
\hline
Threshold frequency $f_0$ & Any frequency works if intense or long enough & No emission at all for $f < f_0$, however intense the light \\
\hline
\rowcolor{popblueL} Time delay & Long delay for weak light & Emission in less than $10^{-9}\ \mathrm{s}$, even for very weak light \\
\hline
$E_{k(\max)}$ vs intensity & $E_{k(\max)}$ increases with intensity & $E_{k(\max)}$ totally independent of intensity \\
\hline
\rowcolor{popblueL} $E_{k(\max)}$ vs frequency & No special role for frequency & $E_{k(\max)}$ increases linearly with $f$ \\
\hline
Effect of intensity & Stronger light $\Rightarrow$ more energetic electrons & Stronger light $\Rightarrow$ \emph{more} electrons, each with the \emph{same} maximum energy \\
\hline
\end{tabularx}
\end{center}

\textbf{The crisis.} Every single prediction fails. The most dramatic failure is the threshold: an intense red spotlight delivers far more energy per second than a faint ultraviolet lamp, yet the red light ejects nothing while the ultraviolet acts instantly. A simple estimate makes the absurdity concrete: treating the electron as absorbing energy from an area of atomic size, classical theory predicts delays of minutes or hours for weak sources --- experiment gives less than a nanosecond. Classical physics, triumphant in explaining interference and diffraction, was bankrupt before this experiment. A radically new model was needed: Einstein's photon hypothesis (developed in the next section).

\begin{attention}[The classic trap: ``brighter red light will eventually work'']
A very common mistake is to believe that increasing the \emph{intensity} of red light (or waiting longer) will eventually cause emission from a metal whose threshold lies in the blue or ultraviolet. \textbf{This is false.} If $f < f_0$, not a single photoelectron is emitted, whatever the intensity and however long you wait. Intensity controls \emph{how many} electrons are emitted per second; frequency alone decides \emph{whether} emission occurs at all and \emph{how energetic} the electrons are.
\end{attention}

\courssub{The Stopping Potential: Measuring the Electron Energy}

\textbf{Idea.} How can we measure the maximum kinetic energy of the photoelectrons? By making them climb an electrostatic ``hill''. If we reverse the supply so that the anode becomes \emph{negative} relative to the cathode ($U_{AC} < 0$), the electric field opposes the motion of the electrons: only the fastest ones reach the anode. As we make the anode more and more negative, the current falls, and at one precise voltage it becomes exactly zero.

\begin{definition}[Stopping potential]
The \cle{stopping potential} $V_s$ (also written $U_0$) is the magnitude of the reverse (retarding) potential difference between anode and cathode for which the photoelectric current just falls to zero: even the most energetic photoelectrons are turned back before reaching the anode.
\end{definition}

\textbf{Derivation (examinable).} Apply the work--energy theorem to the fastest photoelectron, emitted from the cathode with kinetic energy $E_{k(\max)}$ and travelling against a retarding p.d.\ of magnitude $V_s$. The work done by the electric force on the electron (charge $-e$) over the retarding voltage is $W = -eV_s$. At the stopping potential the electron arrives at the anode with zero speed, so
\[
\Delta E_k = W \quad\Longrightarrow\quad 0 - E_{k(\max)} = -eV_s,
\]
hence
\[
\boxed{E_{k(\max)} = e\,V_s}
\]
with $e = 1.60\times 10^{-19}\ \mathrm{C}$. This is why energies in atomic physics are often quoted in \cle{electronvolts}: $1\ \mathrm{eV} = 1.60\times 10^{-19}\ \mathrm{J}$ is the kinetic energy gained by one electron accelerated through $1\ \mathrm{V}$.

\begin{popfigure}
\begin{center}
\begin{tikzpicture}
\begin{axis}[
  width=11cm, height=7.5cm,
  axis lines=middle,
  xlabel={$U_{AC}$ (V)}, ylabel={$I$ ($\mu$A)},
  xmin=-2.6, xmax=4.2, ymin=-0.6, ymax=5.2,
  xtick={-1.5,0,1,2,3,4}, ytick={1,2,3,4},
  every axis x label/.style={at={(ticklabel* cs:1.02)}, anchor=west},
  every axis y label/.style={at={(ticklabel* cs:1.02)}, anchor=south},
]
% High intensity curve
\addplot[very thick, popblue, smooth, domain=-1.5:4, samples=60]
  {4*(1 - exp(-(x+1.5)/0.8))};
% Low intensity curve
\addplot[very thick, poporange, smooth, domain=-1.5:4, samples=60]
  {2*(1 - exp(-(x+1.5)/0.8))};
% Stopping potential marker
\addplot[mark=*, popdark, only marks] coordinates {(-1.5,0)};
\draw[dashed, popdark] (axis cs:-1.5,0) -- (axis cs:-1.5,-0.55);
\node[popdark, font=\small] at (axis cs:-1.5,-0.4) {$-V_s$};
\node[popblue, font=\small, align=center] at (axis cs:2.6,4.35) {high intensity\\ (same $f$)};
\node[poporange, font=\small, align=center] at (axis cs:2.6,1.6) {low intensity\\ (same $f$)};
\node[font=\small, align=center] at (axis cs:1.2,3.1) {saturation\\ current};
\end{axis}
\end{tikzpicture}
\end{center}
\legende{Photoelectric current versus applied voltage for two intensities of the \emph{same} monochromatic light. The stopping potential $V_s$ is the same for both curves (it depends on frequency, not intensity), while the saturation current doubles with intensity.}
\end{popfigure}

The graph summarises the whole section beautifully: both curves start at the \emph{same} stopping potential $-V_s$ (law 4: energy depends on frequency only), but the brighter light gives a larger \emph{saturation current} (law 3: more electrons per second).

\begin{methode}[Using the stopping potential to find $E_{k(\max)}$]
\begin{enumerate}
\item Illuminate the cathode with monochromatic light of known frequency $f$ (check $f > f_0$).
\item Reverse the supply and increase the retarding voltage slowly until the microammeter reads exactly zero.
\item Read the voltmeter: its magnitude is the stopping potential $V_s$.
\item Compute the maximum kinetic energy: $E_{k(\max)} = eV_s$ (in joules), or directly $E_{k(\max)} = V_s$ in electronvolts.
\item Optionally, deduce the maximum speed from $E_{k(\max)} = \tfrac{1}{2}m_e v_{\max}^2$, i.e.\ $v_{\max} = \sqrt{2eV_s/m_e}$.
\end{enumerate}
\end{methode}

\begin{exoresolu}[Stopping potential of a sodium cathode]
Sodium is illuminated with ultraviolet light and emits photoelectrons. The photoelectric current is found to vanish when the retarding potential difference reaches $V_s = 1.85\ \mathrm{V}$.\\
(a) Calculate the maximum kinetic energy of the photoelectrons, in joules and in electronvolts.\\
(b) Calculate the maximum speed of the photoelectrons.\\
(c) The intensity of the ultraviolet source is now doubled. What happens to the stopping potential and to the saturation current?
\tcblower
\textbf{Solution.}\\
(a) At the stopping potential, $E_{k(\max)} = eV_s$:
\[
E_{k(\max)} = 1.60\times 10^{-19} \times 1.85 = 2.96\times 10^{-19}\ \mathrm{J}.
\]
In electronvolts this is simply $E_{k(\max)} = 1.85\ \mathrm{eV}$.\\[4pt]
(b) From $E_{k(\max)} = \tfrac{1}{2}m_e v_{\max}^2$:
\[
v_{\max} = \sqrt{\frac{2E_{k(\max)}}{m_e}} = \sqrt{\frac{2\times 2.96\times 10^{-19}}{9.11\times 10^{-31}}} = \sqrt{6.50\times 10^{11}} \approx 8.06\times 10^{5}\ \mathrm{m\,s^{-1}}.
\]
(c) Doubling the intensity doubles the \emph{number} of photons arriving per second, so the saturation current doubles. But the energy of each photoelectron depends only on the \emph{frequency}, which is unchanged: the stopping potential stays at $V_s = 1.85\ \mathrm{V}$.
\end{exoresolu}

\begin{attention}[Sign conventions with the stopping potential]
The stopping potential is quoted as a \emph{magnitude} (a positive number of volts), but it corresponds to a \emph{reverse} voltage: the anode is negative relative to the cathode, $U_{AC} = -V_s$. In $E_{k(\max)} = eV_s$, use $e = +1.60\times 10^{-19}\ \mathrm{C}$ (the elementary charge, positive). Mixing up the signs is a frequent source of error in the examination.
\end{attention}

\begin{exercice}[Threshold and stopping potential]
A zinc cathode ($\lambda_0 = 370\ \mathrm{nm}$) is illuminated successively by: (i) an intense red laser of wavelength $650\ \mathrm{nm}$; (ii) a very weak ultraviolet lamp of wavelength $250\ \mathrm{nm}$.\\
1. For each source, state whether photoelectrons are emitted, justifying your answer.\\
2. For the case where emission occurs, the stopping potential is measured as $V_s = 1.60\ \mathrm{V}$. Deduce $E_{k(\max)}$ in joules and the maximum speed of the photoelectrons.\\
3. The red laser's power is multiplied by 1000. Does anything change? Comment.
\end{exercice}

\begin{corrige}[Exercise 1]
1. (i) Red light: $f = c/\lambda = 3.00\times 10^{8}/(650\times 10^{-9}) = 4.62\times 10^{14}\ \mathrm{Hz}$, while $f_0 = c/\lambda_0 = 3.00\times 10^{8}/(370\times 10^{-9}) = 8.11\times 10^{14}\ \mathrm{Hz}$. Since $f < f_0$, \textbf{no emission}, despite the high intensity. (ii) Ultraviolet: $f = 3.00\times 10^{8}/(250\times 10^{-9}) = 1.20\times 10^{15}\ \mathrm{Hz} > f_0$: emission occurs, instantly, even though the lamp is weak.\\
2. $E_{k(\max)} = eV_s = 1.60\times 10^{-19}\times 1.60 = 2.56\times 10^{-19}\ \mathrm{J}$; $v_{\max} = \sqrt{2E_{k(\max)}/m_e} = \sqrt{5.62\times 10^{11}} \approx 7.50\times 10^{5}\ \mathrm{m\,s^{-1}}$.\\
3. Nothing changes: the frequency is still below threshold, so no electron is emitted, whatever the power. Only raising the \emph{frequency} above $f_0$ would produce emission.
\end{corrige}

\begin{aretenir}[Key points of this section]
\begin{itemize}
\item The photoelectric effect is the emission of electrons from a metal illuminated by radiation of frequency $f$ greater than a threshold $f_0$ characteristic of the metal.
\item Four experimental laws: threshold frequency; instantaneous emission ($<10^{-9}\ \mathrm{s}$); photoelectric current proportional to intensity; maximum kinetic energy depending on frequency only.
\item Classical wave theory predicts the opposite on every point: this crisis demanded a new model of light.
\item The stopping potential $V_s$ is the reverse voltage that just reduces the current to zero; it measures the maximum kinetic energy through $E_{k(\max)} = eV_s$.
\item Intensity controls \emph{how many} electrons; frequency controls \emph{whether} and \emph{how fast}.
\end{itemize}
\end{aretenir}

\coursec{Photons and Planck's Constant}

In the previous section we saw that the photoelectric effect stubbornly refuses to be explained by the wave theory of light. Three facts stood out like accusations: there is a \cle{threshold frequency} below which no electron is ever ejected, however intense the light; the maximum kinetic energy of the emitted electrons depends only on the \emph{frequency}, not on the intensity; and the emission is \emph{instantaneous}, with no measurable delay even for very weak light. If light were a continuous wave spreading its energy evenly over the metal surface, none of this should happen. A new picture of light is needed. That picture was built in two revolutionary steps: Planck's quantum hypothesis (1900) and Einstein's photon model (1905).

\courssub{Planck's Quantum Hypothesis (1900)}

\textbf{The problem Planck was trying to solve.}\par
Around 1900, physicists were struggling to explain the spectrum of the radiation emitted by a hot body (an idealised emitter called a \emph{black body}). Classical physics predicted that a hot object should emit infinitely more energy at short wavelengths --- the so-called ``ultraviolet catastrophe'' --- in flagrant contradiction with experiment. The German physicist Max Planck found a formula that fitted the experimental curves perfectly, but to derive it he had to make a desperate and, in his own view, almost absurd assumption.

\begin{definition}[Planck's quantum hypothesis]
The energy of the oscillating charges that emit and absorb electromagnetic radiation is not continuous: it can only be exchanged in \cle{discrete packets} called \cle{quanta} (singular: \emph{quantum}). Each quantum of radiation of frequency $f$ carries the energy
\[ E = h f \]
where $h$ is a universal constant of nature called \cle{Planck's constant}.
\end{definition}

In other words, energy at the microscopic scale behaves like money: you can pay 100 FCFA or 200 FCFA, but never 137.4 FCFA if the smallest coin is 100 FCFA. Nature, at the atomic scale, does not allow arbitrary amounts of energy --- only whole multiples of $hf$.

\begin{definition}[Planck's constant]
\cle{Planck's constant} $h$ is the fundamental constant of quantum physics that fixes the size of the energy quantum:
\[ h = 6.63 \times 10^{-34}\ \mathrm{J\,s} \quad (\text{joule-second}). \]
Its unit follows directly from $E = hf$: $[h] = \dfrac{[E]}{[f]} = \dfrac{\mathrm{J}}{\mathrm{Hz}} = \mathrm{J\,s}$.
\end{definition}

\begin{attention}[Why is the quantum world invisible in daily life?]
Planck's constant is \emph{extremely} small: $6.63 \times 10^{-34}\ \mathrm{J\,s}$. A torch beam of frequency $f \approx 5 \times 10^{14}\ \mathrm{Hz}$ has quanta of energy $E = hf \approx 3.3 \times 10^{-19}\ \mathrm{J}$. This is so tiny that the ``graininess'' of light energy is completely undetectable in everyday life, exactly as we do not feel the individual atoms when we touch a smooth table. The quantum structure of energy only reveals itself at the scale of atoms and electrons.
\end{attention}

\courssub{Einstein's Photon Model (1905)}

Planck had quantised the \emph{emission and absorption} of light, but he still believed light itself travelled as a wave. In 1905, Albert Einstein took the idea much further to explain the photoelectric effect.

\begin{definition}[The photon]
Light (and every electromagnetic radiation) consists of a stream of particles of energy called \cle{photons}. Each photon:
\begin{itemize}
\item carries a quantum of energy $E = hf$;
\item travels in vacuum at the speed of light $c = 3.0 \times 10^{8}\ \mathrm{m\,s^{-1}}$;
\item has zero rest mass and carries no electric charge;
\item interacts with matter ``all or nothing'': it is either completely absorbed or not absorbed at all.
\end{itemize}
\end{definition}

\begin{propriete}[Energy of a photon]
The energy of a photon of frequency $f$ and wavelength $\lambda$ (in vacuum) is
\[ E = hf = \frac{hc}{\lambda}, \qquad \text{since } c = f\lambda . \]
With $h = 6.63 \times 10^{-34}\ \mathrm{J\,s}$ and $c = 3.0 \times 10^{8}\ \mathrm{m\,s^{-1}}$, the product $hc = 1.99 \times 10^{-25}\ \mathrm{J\,m}$.
\end{propriete}

This relation shows immediately that \emph{short wavelengths mean energetic photons}: ultraviolet photons are more energetic than visible photons, which are themselves more energetic than infrared photons. This is exactly what the photoelectric experiments demanded.

\begin{popfigure}
\begin{center}
\begin{tikzpicture}[scale=0.95]
% ---- Wave picture (left) ----
\node[font=\bfseries, text=popink] at (2.6,3.4) {Wave picture (classical)};
\draw[popblue, thick, domain=0:4.6, samples=60] plot (\x, {2.2 + 0.45*sin(360*\x/1.15)});
\foreach \y in {1.6,2.0,2.4,2.8} {
  \draw[-latex, popblue!70] (4.7,\y) -- (5.6,\y);
}
\draw[fill=popmuted!40, draw=popink] (5.7,1.2) rectangle (6.3,3.2);
\node[text=popink, font=\small] at (6.0,0.85) {metal};
\node[text=popblue, font=\small, align=center, text width=4.6cm] at (2.6,0.6) {continuous energy spread\\ over the whole front};
% ---- Photon picture (right) ----
\node[font=\bfseries, text=popink] at (10.6,3.4) {Photon picture (Einstein)};
\foreach \x/\y in {7.6/2.0, 8.3/2.6, 9.0/1.8, 9.7/2.4, 10.4/2.9, 8.0/2.9, 9.4/3.0} {
  \fill[poporange] (\x,\y) circle (0.09);
}
\foreach \y in {1.6,2.0,2.4,2.8} {
  \draw[-latex, poporange!80] (10.6,\y) -- (11.5,\y);
}
\draw[fill=popmuted!40, draw=popink] (11.6,1.2) rectangle (12.2,3.2);
\node[text=popink, font=\small] at (11.9,0.85) {metal};
\node[text=poporange!80!black, font=\small, align=center, text width=4.6cm] at (10.6,0.6) {stream of discrete packets,\\ each carrying $E = hf$};
\end{tikzpicture}
\legende{Two pictures of light striking a metal surface. Left: the classical wave spreads its energy continuously over the surface. Right: Einstein's photons arrive as discrete, localised packets of energy $E = hf$.}
\end{center}
\end{popfigure}

\courssub{The Electron-Volt: A Convenient Unit}

Photon energies are of the order of $10^{-19}\ \mathrm{J}$. Writing such numbers is painful and error-prone, so atomic and quantum physicists use a much more convenient unit.

\begin{definition}[The electron-volt]
One \cle{electron-volt} (symbol: $\mathrm{eV}$) is the kinetic energy gained by an electron accelerated from rest through a potential difference of one volt:
\[ 1\ \mathrm{eV} = e \times 1\ \mathrm{V} = 1.60 \times 10^{-19}\ \mathrm{J}. \]
\end{definition}

\begin{methode}[Converting photon energies between joules and electron-volts]
\begin{enumerate}
\item Compute the energy in joules: $E = \dfrac{hc}{\lambda}$ with $\lambda$ \textbf{in metres}.
\item To convert joules $\rightarrow$ electron-volts, \textbf{divide} by $1.60 \times 10^{-19}$:
\[ E\ (\mathrm{eV}) = \frac{E\ (\mathrm{J})}{1.60 \times 10^{-19}\ \mathrm{J/eV}}. \]
\item To convert electron-volts $\rightarrow$ joules, \textbf{multiply} by $1.60 \times 10^{-19}$.
\item \textbf{Quick check:} a visible photon must give an answer of a few eV. If you find $10^{15}$ eV or $10^{-30}$ eV, a unit conversion went wrong.
\end{enumerate}
\end{methode}

\begin{attention}[The classic trap: wavelength not in metres]
In $E = \dfrac{hc}{\lambda}$, the wavelength $\lambda$ \textbf{must} be in metres. GCE questions routinely give $\lambda$ in nanometres. Before any calculation, convert:
\[ 1\ \mathrm{nm} = 10^{-9}\ \mathrm{m}, \qquad 1\ \mathrm{cm} = 10^{-2}\ \mathrm{m}. \]
Forgetting this conversion introduces an error of $10^{9}$ --- your photon becomes absurdly energetic, and the mark is lost.
\end{attention}

\courssub{Photon Energies Across the Spectrum: Worked Examples}

\begin{exoresolu}[Energy of a red photon and a violet photon]
Statement. Calculate, in joules and in electron-volts, the energy of (a) a red photon of wavelength $\lambda_R = 700\ \mathrm{nm}$; (b) a violet photon of wavelength $\lambda_V = 400\ \mathrm{nm}$. Take $h = 6.63 \times 10^{-34}\ \mathrm{J\,s}$, $c = 3.0 \times 10^{8}\ \mathrm{m\,s^{-1}}$, $1\ \mathrm{eV} = 1.60 \times 10^{-19}\ \mathrm{J}$.
\tcblower
\textbf{Solution.}\par
(a) Red photon. Convert the wavelength: $\lambda_R = 700\ \mathrm{nm} = 7.00 \times 10^{-7}\ \mathrm{m}$.
\begin{align*}
E_R &= \frac{hc}{\lambda_R} = \frac{(6.63 \times 10^{-34})(3.0 \times 10^{8})}{7.00 \times 10^{-7}} \\
&= \frac{1.989 \times 10^{-25}}{7.00 \times 10^{-7}} = 2.84 \times 10^{-19}\ \mathrm{J}.
\end{align*}
Converting to electron-volts:
\[ E_R = \frac{2.84 \times 10^{-19}}{1.60 \times 10^{-19}} = 1.8\ \mathrm{eV}. \]
(b) Violet photon. $\lambda_V = 400\ \mathrm{nm} = 4.00 \times 10^{-7}\ \mathrm{m}$.
\[ E_V = \frac{1.989 \times 10^{-25}}{4.00 \times 10^{-7}} = 4.97 \times 10^{-19}\ \mathrm{J} = 3.1\ \mathrm{eV}. \]
\textbf{Conclusion:} visible photons carry energies of a \emph{few electron-volts}, increasing from red ($\approx 1.8\ \mathrm{eV}$) to violet ($\approx 3.1\ \mathrm{eV}$). Violet photons are nearly twice as energetic as red ones --- which is why violet light can eject electrons from metals on which red light has no effect.
\end{exoresolu}

\begin{exemplebox}[Orders of magnitude across the electromagnetic spectrum]
Using $E = hc/\lambda$:
\begin{itemize}
\item Infrared, $\lambda = 1.0 \times 10^{-6}\ \mathrm{m}$: $E \approx 2.0 \times 10^{-19}\ \mathrm{J} \approx 1.2\ \mathrm{eV}$ --- usually too weak to cause photoemission from most metals.
\item Visible, $\lambda \approx 400$--$700\ \mathrm{nm}$: $E \approx 1.8$--$3.1\ \mathrm{eV}$.
\item Ultraviolet, $\lambda = 200\ \mathrm{nm}$: $E \approx 9.9 \times 10^{-19}\ \mathrm{J} \approx 6.2\ \mathrm{eV}$ --- energetic enough for photoemission from many metals.
\end{itemize}
The energy of one photon is fixed by its frequency alone; making the light brighter does \emph{not} make individual photons more energetic.
\end{exemplebox}

\begin{popfigure}
\begin{center}
\begin{tikzpicture}
\begin{axis}[
  ybar, bar width=14pt,
  width=0.86\linewidth, height=5.6cm,
  ymin=0, ymax=3.6,
  ylabel={Photon energy (eV)},
  symbolic x coords={Red, Orange, Yellow, Green, Blue, Violet},
  xtick=data,
  x tick label style={font=\small},
  ytick={0,0.5,1.0,1.5,2.0,2.5,3.0,3.5},
  nodes near coords, nodes near coords style={font=\small},
  every node near coord/.append style={/pgf/number format/fixed, /pgf/number format/precision=2},
  axis lines=left,
  ymajorgrids=true, grid style={popline!50},
]
\addplot[fill=popblue, draw=popink] coordinates {(Red,1.77) (Orange,2.00) (Yellow,2.14) (Green,2.34) (Blue,2.64) (Violet,3.10)};
\end{axis}
\end{tikzpicture}
\legende{Photon energies across the visible spectrum, from red ($\lambda \approx 700\ \mathrm{nm}$, $1.77\ \mathrm{eV}$) to violet ($\lambda \approx 400\ \mathrm{nm}$, $3.10\ \mathrm{eV}$). Energy grows steadily as wavelength decreases.}
\end{center}
\end{popfigure}

\courssub{The Photon Interpretation of Intensity}

If light is a hail of photons, what does its \emph{intensity} mean? In the wave picture, intensity was the energy carried per unit area per unit time by the wave. In the photon picture, the same energy is transported by the photons, so:

\begin{propriete}[Intensity in the photon model]
The intensity $I$ of a monochromatic beam (frequency $f$) equals the number $N$ of photons crossing a unit area per unit time, multiplied by the energy of one photon:
\[ I = N \times hf \quad \Longleftrightarrow \quad N = \frac{I}{hf}. \]
For a source of power $P$ illuminating an area $A$, the number of photons arriving per second on that area is
\[ n = \frac{P}{hf} = \frac{P}{E_{\text{photon}}}. \]
\end{propriete}

\begin{attention}[Exam tip: photons per second from power]
GCE questions very often give the \emph{power} (or intensity and area) of a light source and ask for the number of photons emitted or absorbed per second. The key relation is
\[ \text{number of photons per second} = \frac{\text{power}}{\text{energy of one photon}} = \frac{P}{hf} = \frac{P\lambda}{hc}. \]
Compute $E_{\text{photon}}$ first (in joules!), then divide the power by it.
\end{attention}

\begin{exoresolu}[Counting photons from a lamp]
Statement. A small sodium lamp emits yellow light of wavelength $\lambda = 589\ \mathrm{nm}$ with a useful light power $P = 2.0\ \mathrm{W}$. Calculate the number of photons emitted per second. ($h = 6.63 \times 10^{-34}\ \mathrm{J\,s}$, $c = 3.0 \times 10^{8}\ \mathrm{m\,s^{-1}}$.)
\tcblower
\textbf{Solution.}\par
Energy of one photon (with $\lambda = 589 \times 10^{-9}\ \mathrm{m}$):
\[ E_{\text{photon}} = \frac{hc}{\lambda} = \frac{(6.63 \times 10^{-34})(3.0 \times 10^{8})}{589 \times 10^{-9}} = 3.38 \times 10^{-19}\ \mathrm{J}. \]
Number of photons per second:
\[ n = \frac{P}{E_{\text{photon}}} = \frac{2.0}{3.38 \times 10^{-19}} \approx 5.9 \times 10^{18}\ \text{photons per second}. \]
Even a modest lamp emits nearly six billion billion photons every second --- which is why light usually \emph{looks} continuous.
\end{exoresolu}

\courssub{How the Photon Model Explains the Photoelectric Effect}

We can now close the loop and see how naturally the photon model dissolves the three paradoxes of the previous section.

\textbf{1. The threshold frequency.} In the photon picture, an electron is ejected in a \emph{one photon -- one electron} collision: a single photon is absorbed completely by a single electron. The electron can escape the metal only if the photon's energy $hf$ is at least equal to the minimum work needed to extract it, the \cle{work function} $W_0$ of the metal (studied in detail in the next section):
\[ hf \geq W_0 \quad \Longleftrightarrow \quad f \geq f_0 = \frac{W_0}{h}. \]
If $f < f_0$, each photon is individually too weak. Increasing the intensity merely sends \emph{more} insufficient photons --- like throwing millions of ping-pong balls at a door: none can break it, however many you throw. One energetic ``bullet'' (a photon with $f \geq f_0$) succeeds where an army of weak ones fails.

\textbf{2. Instantaneous emission.} The energy is delivered in a single, localised packet. As soon as a photon with $f \geq f_0$ hits a suitable electron, the electron is ejected immediately --- there is no need to accumulate energy gradually from a spread-out wave front. Even at very low intensity, the \emph{first} photon that arrives can eject an electron, with no measurable delay (less than about $10^{-9}\ \mathrm{s}$ experimentally).

\textbf{3. Kinetic energy depends on frequency, not intensity.} Since the interaction is one photon for one electron, the surplus energy $hf - W_0$ becomes the electron's kinetic energy: it depends on $f$ alone. Raising the intensity at fixed $f$ increases the \emph{number} of photons, hence the \emph{number} of ejected electrons (the current), but not their individual energies.

\begin{exoresolu}[Threshold frequency of sodium]
Statement. The work function of sodium is $W_0 = 2.3\ \mathrm{eV}$. (a) Calculate the threshold frequency $f_0$ and the corresponding threshold wavelength $\lambda_0$ of sodium. (b) State, with justification, whether green light of wavelength $550\ \mathrm{nm}$ can eject electrons from sodium. ($h = 6.63 \times 10^{-34}\ \mathrm{J\,s}$, $c = 3.0 \times 10^{8}\ \mathrm{m\,s^{-1}}$, $1\ \mathrm{eV} = 1.60 \times 10^{-19}\ \mathrm{J}$.)
\tcblower
\textbf{Solution.}\par
(a) Convert the work function to joules:
\[ W_0 = 2.3 \times 1.60 \times 10^{-19} = 3.68 \times 10^{-19}\ \mathrm{J}. \]
Threshold frequency:
\[ f_0 = \frac{W_0}{h} = \frac{3.68 \times 10^{-19}}{6.63 \times 10^{-34}} = 5.6 \times 10^{14}\ \mathrm{Hz}. \]
Threshold wavelength:
\[ \lambda_0 = \frac{c}{f_0} = \frac{3.0 \times 10^{8}}{5.6 \times 10^{14}} = 5.4 \times 10^{-7}\ \mathrm{m} = 540\ \mathrm{nm}. \]
(b) Green light has $\lambda = 550\ \mathrm{nm} > \lambda_0 = 540\ \mathrm{nm}$, i.e.\ $f < f_0$. Its photon energy is
\[ E = \frac{hc}{\lambda} = \frac{1.989 \times 10^{-25}}{5.50 \times 10^{-7}} = 3.62 \times 10^{-19}\ \mathrm{J} = 2.3\ \mathrm{eV}, \]
marginally \emph{below} the work function. Therefore green light of $550\ \mathrm{nm}$ cannot eject electrons from sodium, however intense it may be.
\end{exoresolu}

\begin{aretenir}[Key points of this section]
\begin{itemize}
\item Planck (1900): energy is exchanged in quanta $E = hf$, with $h = 6.63 \times 10^{-34}\ \mathrm{J\,s}$.
\item Einstein (1905): light itself is a stream of photons, each of energy $E = hf = \dfrac{hc}{\lambda}$, travelling at speed $c$.
\item $1\ \mathrm{eV} = 1.60 \times 10^{-19}\ \mathrm{J}$; visible photons carry a few eV.
\item Intensity $=$ (photons per unit area per unit time) $\times$ (photon energy); photons per second $= P/(hf)$.
\item The one photon -- one electron picture explains the threshold frequency $f_0 = W_0/h$ and the instantaneous character of photoemission.
\end{itemize}
\end{aretenir}

\begin{exercice}[Photons from a UV source]
An ultraviolet source emits light of wavelength $\lambda = 250\ \mathrm{nm}$ with a power of $0.50\ \mathrm{W}$.
\begin{enumerate}
\item Calculate the energy of one photon, in joules and in electron-volts.
\item How many photons does the source emit per second?
\item The work function of zinc is $W_0 = 4.3\ \mathrm{eV}$. Can this radiation eject electrons from a zinc plate? Justify.
\end{enumerate}
\end{exercice}

\begin{corrige}[Exercise 1]
1. $\lambda = 250 \times 10^{-9}\ \mathrm{m}$.
\[ E = \frac{hc}{\lambda} = \frac{1.989 \times 10^{-25}}{2.50 \times 10^{-7}} = 7.96 \times 10^{-19}\ \mathrm{J} = \frac{7.96 \times 10^{-19}}{1.60 \times 10^{-19}} = 5.0\ \mathrm{eV}. \]
2. $n = \dfrac{P}{E} = \dfrac{0.50}{7.96 \times 10^{-19}} \approx 6.3 \times 10^{17}$ photons per second.
3. Yes: $E = 5.0\ \mathrm{eV} > W_0 = 4.3\ \mathrm{eV}$, so each photon has enough energy to extract an electron from zinc (the surplus $0.7\ \mathrm{eV}$ becomes kinetic energy).
\end{corrige}

\begin{savaistu}[A Nobel Prize for the photon]
It was for his explanation of the photoelectric effect --- and \emph{not} for relativity --- that Albert Einstein received the Nobel Prize in Physics in 1921. The idea that light comes in particles was so shocking that even Planck initially opposed it. Today, photons are the workhorses of modern technology: solar panels powering homes across Cameroon, photodiodes in fibre-optic communications, and the sensors of every smartphone camera all rely on photon--electron interactions.
\end{savaistu}

\coursec{Einstein's Photoelectric Equation}

In the previous section we established that light behaves as a stream of energy quanta called \cle{photons}, each carrying an energy $E = hf = \dfrac{hc}{\lambda}$, where $h = \SI{6.63e-34}{J.s}$ is Planck's constant. We now use this revolutionary idea to do what the wave theory could never do: explain the photoelectric effect \emph{quantitatively}. Einstein's photoelectric equation, published in 1905, is nothing more than the conservation of energy applied to the interaction between one photon and one electron in a metal. It is one of the most frequently examined relations at the GCE Advanced Level, so we shall build it carefully, term by term.

\courssub{The Work Function and the Threshold Frequency}

\textbf{Why an electron needs energy to escape.}\par

\textbf{Intuition first.} Imagine a child at the bottom of a deep well in the village. To climb out, the child must spend energy working against gravity; the deeper the well, the more energy is needed. An electron inside a metal is in a similar situation: it is attracted by the positive ions of the metallic lattice, so it sits in a kind of ``potential well''. To tear it away from the metal surface, energy must be supplied to it. Different metals hold their electrons with different strengths, just as wells have different depths.

\begin{definition}[Work function]
The \cle{work function} $W_0$ of a metal is the \textbf{minimum energy} required to extract an electron from the surface of that metal. It is a property of the metal alone (not of the light used). It is related to the \cle{threshold frequency} $f_0$ and the \cle{threshold wavelength} $\lambda_0$ by
\[
W_0 = h f_0 = \frac{hc}{\lambda_0}.
\]
\end{definition}

\begin{definition}[Threshold frequency and threshold wavelength]
The \cle{threshold frequency} $f_0$ is the \textbf{lowest frequency} of radiation that can eject electrons from a given metal. The \cle{threshold wavelength} $\lambda_0 = c/f_0$ is the \textbf{longest wavelength} that can do so. Radiation with $f < f_0$ (equivalently $\lambda > \lambda_0$) can never cause photoemission, however intense it may be.
\end{definition}

\begin{exemplebox}[Typical values for common metals]
The table below gives approximate values frequently used in GCE problems. Work functions are often quoted in electronvolts: $1\ \mathrm{eV} = 1.60\times 10^{-19}\ \mathrm{J}$.
\begin{center}
\begin{tabularx}{\linewidth}{|l|X|X|X|}
\hline
\rowcolor{popblueL} \textbf{Metal} & \textbf{$W_0$ (eV)} & \textbf{$f_0$ (Hz)} & \textbf{$\lambda_0$ (nm)} \\
\hline
Sodium (Na) & $2.3$ & $5.6\times 10^{14}$ & $540$ \\
\hline
Calcium (Ca) & $2.9$ & $7.0\times 10^{14}$ & $430$ \\
\hline
Zinc (Zn) & $4.3$ & $1.0\times 10^{15}$ & $290$ \\
\hline
\end{tabularx}
\end{center}
Notice that sodium responds to visible light (green), whereas zinc needs ultraviolet radiation: this is why zinc plates are used in classic demonstrations with UV lamps.
\end{exemplebox}

\begin{attention}[The work function is a MINIMUM]
$W_0$ is the energy needed to extract the \emph{least tightly bound} electrons, those right at the surface. Electrons lying deeper inside the metal need \emph{more} energy to escape. This is why, as we shall see, the kinetic energy computed from Einstein's equation is a \emph{maximum}.
\end{attention}

\courssub{Statement and Derivation of Einstein's Equation}

\textbf{One photon, one electron.}\par

Einstein's key assumption is that the interaction is \textbf{individual}: a single photon is completely absorbed by a single electron. The photon ceases to exist and hands over \emph{all} of its energy $hf$ to the electron in one indivisible transaction. There is no gradual accumulation of energy, which immediately explains why emission is instantaneous and why weak light below the threshold produces nothing.

The electron then spends part of this energy to escape from the metal (at least $W_0$), and whatever remains appears as kinetic energy.

\begin{propriete}[Einstein's photoelectric equation]
For monochromatic radiation of frequency $f$ incident on a metal of work function $W_0$:
\[
\boxed{\,hf = W_0 + E_{k(\max)}\,}
\]
where $E_{k(\max)} = \dfrac{1}{2}m_e v_{\max}^2$ is the \textbf{maximum} kinetic energy of the emitted photoelectrons. Equivalently, using the stopping potential $V_s$ (the reverse voltage that just stops the most energetic electrons), $E_{k(\max)} = eV_s$, so
\[
hf = W_0 + eV_s.
\]
This equation is a direct statement of the \textbf{conservation of energy}. The derivation (energy balance) is examinable and you must be able to reproduce it in words.
\end{propriete}

\textbf{Derivation (energy balance).}\par
\begin{itemize}
\item \textbf{Energy received:} the electron absorbs one photon of energy $hf$.
\item \textbf{Energy spent:} to leave the metal, the electron must do work at least equal to $W_0$ against the attractive forces of the lattice.
\item \textbf{Energy remaining:} by conservation of energy, the surplus is carried away as kinetic energy: $E_k = hf - W_0$. For a surface electron (minimum possible expenditure $W_0$), this kinetic energy is maximal, hence $E_{k(\max)} = hf - W_0$.
\end{itemize}

\begin{popfigure}
\begin{center}
\begin{tikzpicture}[scale=1.0, font=\small]
% Metal block
\fill[popblueL] (0,-1.6) rectangle (7.5,0);
\draw[thick, popdark] (0,0) -- (7.5,0);
\node[popmuted] at (6.6,-1.25) {metal};
% Photon arrow
\draw[very thick, poporange, decorate, decoration={snake, amplitude=1.2mm, segment length=4mm}] (0.3,1.6) -- (2.3,0.35);
\node[poporange, align=center] at (1.0,2.0) {photon $hf$};
% Electron escaping
\draw[very thick, popgreen, ->] (2.6,-0.25) .. controls (3.6,0.9) .. (5.0,1.3);
\fill[popgreen] (2.6,-0.25) circle (2pt);
\node[popgreen] at (5.6,1.35) {$e^-$: $E_{k(\max)}$};
% Work function barrier
\draw[very thick, poppink, ->] (3.4,-1.2) -- (3.4,-0.05);
\node[poppink, align=center, text width=2.6cm] at (4.6,-0.9) {escape cost $W_0 = hf_0$};
% Balance
\node[popdark, align=center] at (3.75,-2.3) {$\underbrace{hf}_{\text{received}} = \underbrace{W_0}_{\text{spent escaping}} + \underbrace{E_{k(\max)}}_{\text{left over}}$};
\end{tikzpicture}
\end{center}
\legende{Energy balance for one absorbed photon: its energy is split between the work needed to escape the metal and the kinetic energy of the emitted electron.}
\end{popfigure}

\courssub{Condition for Photoemission}

Since kinetic energy can never be negative, emission is only possible if
\[
E_{k(\max)} = hf - W_0 \geq 0
\quad\Longleftrightarrow\quad
hf \geq W_0
\quad\Longleftrightarrow\quad
\boxed{\,f \geq f_0 \quad\text{or equivalently}\quad \lambda \leq \lambda_0.\,}
\]

At the exact threshold ($f = f_0$), the electron is ejected with \emph{zero} kinetic energy: it just barely escapes. Below the threshold, each photon simply does not carry enough energy, and because absorption is all-or-nothing, no electron is ever emitted --- even if billions of photons arrive every second.

\begin{attention}[Intensity versus frequency --- do not mix them up!]
This is the classic GCE trap:
\begin{itemize}
\item Increasing the \textbf{intensity} of the light (with $f > f_0$) means \emph{more photons per second}, hence \emph{more electrons emitted per second} (a larger photoelectric current). The kinetic energy of each electron is \textbf{unchanged}.
\item Increasing the \textbf{frequency} means \emph{more energetic photons}, hence \emph{more energetic electrons} (larger $E_{k(\max)}$ and larger stopping potential). The number of electrons is essentially unchanged.
\end{itemize}
Intensity controls the \emph{number}; frequency controls the \emph{energy}. If $f < f_0$, no intensity, however large, can produce emission.
\end{attention}

\begin{attention}[$E_{k(\max)}$ is a MAXIMUM]
Einstein's equation gives the kinetic energy of the \emph{most energetic} electrons only --- those ejected from the very surface. Electrons from deeper layers lose extra energy in collisions on their way out and emerge with \emph{less} kinetic energy, anywhere between $0$ and $E_{k(\max)}$. Never write ``all photoelectrons have kinetic energy $hf - W_0$''.
\end{attention}

\courssub{The Graph of $E_{k(\max)}$ against $f$}

Rewriting the equation as $E_{k(\max)} = hf - W_0$, we recognise the equation of a straight line $y = mx + c$ with:
\begin{itemize}
\item \textbf{gradient} $= h$ (Planck's constant) --- the same for \emph{all} metals;
\item \textbf{intercept on the frequency axis} $= f_0$ (the threshold frequency, where $E_{k(\max)} = 0$);
\item \textbf{intercept on the energy axis} $= -W_0$.
\end{itemize}

\begin{popfigure}
\begin{center}
\begin{tikzpicture}
\begin{axis}[
    width=11cm, height=7.5cm,
    axis lines=left,
    xlabel={$f$ / $10^{14}$ Hz}, ylabel={$E_{k(\max)}$ / $10^{-19}$ J},
    xmin=0, xmax=12, ymin=-4.5, ymax=5,
    xtick={5.6}, xticklabels={$f_0$},
    ytick={-3.68}, yticklabels={$-W_0$},
    domain=0:12, samples=50, thick,
    clip=false
]
\addplot[popblue, very thick, domain=0:12] {0.663*x - 3.68};
\addplot[popmuted, dashed] coordinates {(5.6,0) (5.6,-3.68) (0,-3.68)};
% gradient triangle
\addplot[poporange, thick] coordinates {(8,1.624) (11,1.624) (11,3.613)};
\node[poporange] at (9.5,1.15) {$\Delta f$};
\node[poporange] at (11.9,2.6) {$\Delta E_k$};
\node[popdark, align=center] at (8.6,4.4) {gradient $= \dfrac{\Delta E_k}{\Delta f} = h$};
\node[popblue] at (3.2,3.4) {$E_{k(\max)} = hf - W_0$};
\end{axis}
\end{tikzpicture}
\end{center}
\legende{Maximum kinetic energy versus frequency (here for sodium). The line cuts the frequency axis at $f_0$ and the energy axis at $-W_0$; its gradient is Planck's constant $h$.}
\end{popfigure}

\begin{methode}[Extracting $h$, $f_0$ and $W_0$ from the graph]
\begin{enumerate}
\item Plot $E_{k(\max)}$ (in joules) on the vertical axis against $f$ (in hertz) on the horizontal axis and draw the best straight line.
\item \textbf{Gradient:} pick two well-separated points on the line and compute $h = \dfrac{\Delta E_{k(\max)}}{\Delta f}$. Watch the powers of ten on the axes!
\item \textbf{Threshold frequency:} read $f_0$ where the line crosses the $f$-axis ($E_{k(\max)} = 0$).
\item \textbf{Work function:} either read the magnitude of the vertical intercept, $W_0 = |{-W_0}|$, or compute $W_0 = hf_0$. Convert to eV by dividing by $1.60\times 10^{-19}$ if required.
\item \textbf{Threshold wavelength:} $\lambda_0 = c/f_0$.
\end{enumerate}
\end{methode}

\courssub{Link with the Stopping Potential}

In the photocell experiment, we apply a reverse (retarding) voltage between the emitter and the collector and increase it until even the fastest electrons are turned back: the current falls to zero. The work done by the retarding field on an electron of charge $e$ over the \cle{stopping potential} $V_s$ equals the maximum kinetic energy:
\[
eV_s = E_{k(\max)} = hf - W_0
\qquad\Longrightarrow\qquad
V_s = \frac{h}{e}f - \frac{W_0}{e}.
\]
A graph of $V_s$ against $f$ is therefore also a straight line, of gradient $h/e$, with the same frequency intercept $f_0$. This was the basis of Millikan's famous experiments which measured $h$ and confirmed Einstein's theory.

\begin{exoresolu}[Photoelectrons from a sodium surface]
\textbf{Statement.} Light of wavelength $\lambda = \SI{400}{nm}$ falls on a sodium surface of work function $W_0 = \SI{2.28}{eV}$. Take $h = \SI{6.63e-34}{J.s}$, $c = \SI{3.00e8}{m/s}$, $e = \SI{1.60e-19}{C}$. Determine: (a) the threshold frequency and threshold wavelength of sodium; (b) the maximum kinetic energy of the photoelectrons, in joules and in eV; (c) the stopping potential.
\tcblower
\textbf{Solution.}
(a) Convert the work function to joules:
\[
W_0 = 2.28 \times 1.60\times 10^{-19} = 3.65\times 10^{-19}\ \mathrm{J}.
\]
Threshold frequency:
\[
f_0 = \frac{W_0}{h} = \frac{3.65\times 10^{-19}}{6.63\times 10^{-34}} = 5.50\times 10^{14}\ \mathrm{Hz}.
\]
Threshold wavelength:
\[
\lambda_0 = \frac{c}{f_0} = \frac{3.00\times 10^{8}}{5.50\times 10^{14}} = 5.45\times 10^{-7}\ \mathrm{m} = 545\ \mathrm{nm}.
\]
Since $\lambda = 400\ \mathrm{nm} < \lambda_0$, emission does occur.

(b) Photon energy:
\[
hf = \frac{hc}{\lambda} = \frac{6.63\times 10^{-34}\times 3.00\times 10^{8}}{4.00\times 10^{-7}} = 4.97\times 10^{-19}\ \mathrm{J}.
\]
Einstein's equation:
\[
E_{k(\max)} = hf - W_0 = 4.97\times 10^{-19} - 3.65\times 10^{-19} = 1.32\times 10^{-19}\ \mathrm{J} = 0.83\ \mathrm{eV}.
\]

(c) Stopping potential:
\[
V_s = \frac{E_{k(\max)}}{e} = \frac{1.32\times 10^{-19}}{1.60\times 10^{-19}} = 0.83\ \mathrm{V}.
\]
Note the shortcut: when $E_{k(\max)}$ is expressed in eV, the stopping potential in volts has the \emph{same numerical value}.
\end{exoresolu}

\begin{exoresolu}[Determining Planck's constant from experimental data]
\textbf{Statement.} In a photoelectric experiment with a calcium surface, the stopping potentials measured for two frequencies are: $V_{s1} = \SI{0.50}{V}$ at $f_1 = \SI{8.0e14}{Hz}$ and $V_{s2} = \SI{1.60}{V}$ at $f_2 = \SI{1.1e15}{Hz}$. Deduce (a) Planck's constant, (b) the work function and the threshold frequency of calcium.
\tcblower
\textbf{Solution.}
(a) From $eV_s = hf - W_0$, subtracting the two measurements eliminates $W_0$:
\[
e(V_{s2} - V_{s1}) = h(f_2 - f_1)
\quad\Longrightarrow\quad
h = \frac{e(V_{s2}-V_{s1})}{f_2 - f_1} = \frac{1.60\times 10^{-19}\times (1.60 - 0.50)}{1.1\times 10^{15} - 8.0\times 10^{14}} = \frac{1.76\times 10^{-19}}{3.0\times 10^{14}} = 5.9\times 10^{-34}\ \mathrm{J\,s}.
\]
(b) Using the first measurement:
\[
W_0 = hf_1 - eV_{s1} = (5.9\times 10^{-34}\times 8.0\times 10^{14}) - (1.60\times 10^{-19}\times 0.50) = 4.7\times 10^{-19} - 0.80\times 10^{-19} = 3.9\times 10^{-19}\ \mathrm{J} \approx 2.5\ \mathrm{eV},
\]
and
\[
f_0 = \frac{W_0}{h} = \frac{3.9\times 10^{-19}}{5.9\times 10^{-34}} \approx 6.7\times 10^{14}\ \mathrm{Hz}.
\]
(Small discrepancies with accepted values come from the rounded data --- a good lesson in experimental uncertainty.)
\end{exoresolu}

\begin{exercice}[GCE practice]
Ultraviolet radiation of frequency $\SI{1.2e15}{Hz}$ is incident on a zinc plate of work function $\SI{4.3}{eV}$.
\begin{enumerate}
\item Show that photoemission occurs and calculate the threshold wavelength of zinc.
\item Calculate the maximum kinetic energy of the emitted electrons in eV.
\item Calculate the stopping potential and the maximum speed of the photoelectrons ($m_e = \SI{9.11e-31}{kg}$).
\item The intensity of the radiation is now doubled, the frequency being unchanged. State what happens to the answers of parts 2 and 3 and to the photoelectric current.
\end{enumerate}
\end{exercice}

\begin{corrige}[Exercise 1]
1. $W_0 = 4.3\times 1.60\times 10^{-19} = 6.88\times 10^{-19}\ \mathrm{J}$; $f_0 = W_0/h = 1.04\times 10^{15}\ \mathrm{Hz} < 1.2\times 10^{15}\ \mathrm{Hz}$, so emission occurs. $\lambda_0 = c/f_0 = 2.9\times 10^{-7}\ \mathrm{m} = 290\ \mathrm{nm}$.

2. $hf = 6.63\times 10^{-34}\times 1.2\times 10^{15} = 7.96\times 10^{-19}\ \mathrm{J} = 4.97\ \mathrm{eV}$; $E_{k(\max)} = 4.97 - 4.3 = 0.67\ \mathrm{eV}$.

3. $V_s = 0.67\ \mathrm{V}$; $E_{k(\max)} = 0.67\times 1.60\times 10^{-19} = 1.07\times 10^{-19}\ \mathrm{J}$, so
\[
v_{\max} = \sqrt{\frac{2E_{k(\max)}}{m_e}} = \sqrt{\frac{2\times 1.07\times 10^{-19}}{9.11\times 10^{-31}}} \approx 4.9\times 10^{5}\ \mathrm{m\,s^{-1}}.
\]

4. $E_{k(\max)}$, $V_s$ and $v_{\max}$ are \emph{unchanged} (frequency is unchanged); the photoelectric current \emph{doubles} because twice as many photons eject twice as many electrons per second.
\end{corrige}

\begin{aretenir}[Key points]
\begin{itemize}
\item Work function: $W_0 = hf_0 = hc/\lambda_0$ --- minimum energy to extract an electron; a property of the metal.
\item Einstein's equation is conservation of energy: $hf = W_0 + E_{k(\max)} = W_0 + eV_s$.
\item Emission requires $f \geq f_0$ ($\lambda \leq \lambda_0$); at the threshold the electron escapes with zero kinetic energy.
\item The $E_{k(\max)}$--$f$ graph is a straight line of gradient $h$, frequency intercept $f_0$ and energy intercept $-W_0$; the $V_s$--$f$ graph has gradient $h/e$.
\item Frequency controls the \emph{energy} of the electrons; intensity controls their \emph{number}.
\end{itemize}
\end{aretenir}

\coursec{Wave-Particle Duality and Applications}

In the previous section we established Einstein's photoelectric equation, $hf = W_0 + E_{k,\max}$, and saw that it explains every feature of the photoelectric effect \emph{provided} we accept that light arrives in discrete packets of energy --- photons. Yet in your earlier studies of optics you proved, just as convincingly, that light undergoes interference and diffraction, phenomena that only \emph{waves} can produce. Are we facing a contradiction? No. We are facing one of the deepest truths of modern physics, and the subject of this final section: the \cle{wave-particle duality} of light.

\courssub{Two Bodies of Evidence}

\textbf{Evidence for the wave nature of light.}\par
Three families of experiments force a wave description upon us:

\begin{itemize}
\item \textbf{Interference.} When light from a single source passes through two narrow, close slits (Young's double-slit experiment), the screen shows alternate bright and dark fringes. Bright fringes occur where two waves arrive \emph{in phase} and reinforce; dark fringes where they arrive \emph{out of phase} and cancel. Superposition of this kind is meaningless for classical particles.
\item \textbf{Diffraction.} Light spreads out after passing through a narrow slit or around the edge of an obstacle, producing fringes whose spacing depends on the wavelength $\lambda$ and the aperture width. A stream of classical particles would simply cast a sharp shadow.
\item \textbf{Polarisation.} Light can be polarised by a Polaroid filter or by reflection, showing that its vibrations are \emph{transverse}. Only transverse waves can be polarised.
\end{itemize}

\textbf{Evidence for the particle nature of light.}\par
Three phenomena resist any wave explanation:

\begin{itemize}
\item \textbf{The photoelectric effect} (this chapter): instantaneous emission, a threshold frequency $f_0$ independent of intensity, and $E_{k,\max}$ increasing linearly with $f$. All are explained by photons of energy $E = hf$, none by a continuous wave.
\item \textbf{Black-body radiation.} The spectrum of radiation emitted by a hot body cannot be derived from classical wave theory (which predicts the ``ultraviolet catastrophe''). Planck succeeded only by assuming energy is exchanged in quanta $hf$.
\item \textbf{The Compton effect} (mention only): when X-rays scatter off electrons, the scattered radiation has a \emph{longer} wavelength, exactly as predicted if a photon collides with an electron like a billiard ball, conserving energy and momentum. This shows the photon even carries momentum $p = h/\lambda$.
\end{itemize}

\begin{popfigure}
\begin{tikzpicture}[
  exp/.style={rectangle, rounded corners, draw=popblue, fill=popblueL, text width=3.1cm, align=center, font=\small, minimum height=0.9cm},
  mod/.style={rectangle, rounded corners, draw=popdark, fill=popgreenL, text width=3.4cm, align=center, font=\small\bfseries, minimum height=1cm},
  arr/.style={-{Stealth[length=2.5mm]}, thick, popmuted}]
\node[mod, align=center] (wave) at (0,2.6){WAVE model\\ $f$, $\lambda$, amplitude};
\node[mod, align=center] (part) at (0,-2.6){PARTICLE model\\ photon: $E=hf$, $p=h/\lambda$};
\node[exp] (inter) at (-5,2.6) {Interference (Young's slits)};
\node[exp] (diffr) at (-5,0.9) {Diffraction, polarisation};
\node[exp] (photo) at (5,2.6) {Photoelectric effect};
\node[exp] (comp) at (5,0.9) {Black body, Compton effect};
\draw[arr] (inter) -- (wave);
\draw[arr] (diffr) -- (wave);
\draw[arr] (photo) -- (part);
\draw[arr] (comp) -- (part);
\node[rectangle, draw=poporange, thick, fill=white, align=center, font=\small\bfseries] (dual) at (0,0) {LIGHT: one entity,\\ two aspects};
\draw[arr, poporange] (wave) -- (dual);
\draw[arr, poporange] (part) -- (dual);
\end{tikzpicture}
\legende{Concept map: each experiment reveals one aspect of light. The two models are not rivals; both are faces of a single reality.}
\end{popfigure}

\courssub{Statement of Wave-Particle Duality}

\begin{propriete}[Wave-particle duality of light]
Light exhibits \textbf{both} wave-like and particle-like properties. It propagates through space as a wave (undergoing interference, diffraction and polarisation) but exchanges energy and momentum with matter in discrete quanta called \cle{photons} (photoelectric effect, Compton scattering, black-body radiation). The aspect revealed in any observation \emph{depends on the experiment performed}; no single classical picture describes all phenomena.
\end{propriete}

\begin{definition}[Dual nature of radiation]
Every quantum phenomenon associates a \textbf{wave} (characterised by frequency $f$ and wavelength $\lambda$) with a \textbf{particle} (characterised by energy $E$ and momentum $p$). The two descriptions are linked by Planck's constant:
\[ E = hf, \qquad p = \frac{h}{\lambda}, \qquad \text{with } \lambda f = c. \]
\end{definition}

The photon is precisely the \emph{bridge} between the two aspects: the relation $E = hf$ places on the left a \emph{particle} quantity (the energy of one photon) and on the right a \emph{wave} quantity (the frequency of the radiation). Planck's constant $h = \SI{6.63e-34}{J.s}$ is the conversion factor between the two languages.

\begin{methode}[Which model should I use?]
\begin{enumerate}
\item Look at what the phenomenon involves.
\item \textbf{Propagation} of light through space or apparatus (interference fringes, diffraction patterns, polarisation, refraction): use the \textbf{wave model}.
\item \textbf{Interaction or detection} --- energy exchanged with matter in localised events (photoelectric emission, absorption by a photocell, scintillation, Compton scattering): use the \textbf{photon model}.
\item If an essay asks about \emph{both}, present one experiment per aspect and conclude with the duality statement.
\end{enumerate}
\end{methode}

\begin{attention}[The classic mistake]
Do \textbf{not} write that light is ``sometimes a wave and sometimes a particle'', as if it changed its identity from Monday to Tuesday. Light is a \emph{single} physical entity whose complete behaviour cannot be reduced to either classical picture. The wave and the particle are two complementary \emph{models} --- two partial descriptions --- and the experiment you perform selects which aspect becomes visible. Examiners penalise the ``sometimes... sometimes'' phrasing.
\end{attention}

\courssub{Extension: Matter Waves (de Broglie)}

\begin{savaistu}[If waves are particles, are particles waves?]
In 1924 Louis de Broglie proposed the symmetrical idea: \emph{matter} also has wave properties. Every particle of momentum $p$ is associated with a wavelength
\[ \lambda = \frac{h}{p} = \frac{h}{mv}. \]
The hypothesis was confirmed when electrons were observed to diffract through crystals, exactly like X-rays. This is the principle of the \textbf{electron microscope}, which resolves details far smaller than any optical microscope because electron wavelengths can be extremely short. (This extension goes beyond the GCE syllabus but is a favourite ``further thinking'' remark.)
\end{savaistu}

\courssub{Technological Applications of the Photoelectric Effect}

The photoelectric effect is not an examination curiosity: it powers devices you meet daily in Yaoundé, Douala or Bamenda.

\begin{itemize}
\item \textbf{Photocells and automatic light switches.} A photosensitive surface emits electrons when illuminated; the resulting current (or its absence at dusk) operates street lamps, automatic doors and burglar alarms.
\item \textbf{Photovoltaic solar panels.} Photons absorbed in a semiconductor junction liberate charge carriers, producing a direct voltage. Each absorbed photon of energy $hf \geq E_{\text{gap}}$ frees one electron--hole pair: the photon picture directly explains why panel output depends on the \emph{number} of suitable photons, and why low-energy photons are useless whatever the brightness.
\item \textbf{Photodiodes and image sensors.} Digital cameras and phone cameras contain millions of tiny photodiodes; each converts incoming photons into a measurable charge, building the picture pixel by pixel.
\item \textbf{Photomultipliers.} A single photon ejects one electron, which is accelerated onto a series of electrodes (dynodes), each releasing several more electrons. The cascade turns one photon into a measurable pulse of millions of electrons --- used in scanners and radiation detectors.
\end{itemize}

\begin{popfigure}
\begin{tikzpicture}[scale=1.0]
% sun and photons
\shade[inner color=popgold, outer color=white] (-4.6,2.6) circle (0.55);
\node[font=\small] at (-4.6,3.5) {Sun};
\foreach \y in {2.9,2.3,1.7}{
  \draw[-{Stealth[length=2.5mm]}, very thick, poporange] (-3.9,\y) -- (-1.6,\y);
}
\node[font=\small, poporange] at (-2.7,3.25) {photons $hf$};
% cell layers
\fill[popblueL] (-1.5,0.4) rectangle (2.5,1.4);
\fill[popgreenL] (-1.5,-0.6) rectangle (2.5,0.4);
\draw[popdark, thick] (-1.5,-0.6) rectangle (2.5,1.4);
\node[font=\small] at (0.5,0.95) {n-type layer};
\node[font=\small] at (0.5,-0.1) {p-type layer};
\draw[popdark, dashed] (-1.5,0.4) -- (2.5,0.4);
\node[font=\scriptsize, popdark] at (0.5,0.62) {junction};
% electron movement
\draw[-{Stealth[length=2.5mm]}, thick, poppurple] (0.2,-0.3) -- (0.2,1.1);
\node[font=\scriptsize, poppurple, right] at (0.3,0.2) {$e^-$ freed};
% external circuit
\draw[thick, popdark] (2.5,1.1) -- (3.6,1.1) -- (3.6,-0.3) -- (2.5,-0.3);
\draw[thick, popdark] (3.6,0.15) circle (0.28);
\node[font=\small] at (3.6,0.15) {L};
\node[font=\scriptsize, below] at (3.6,-0.4) {load};
\draw[-{Stealth[length=2mm]}, thick, poppurple] (2.6,1.1) -- (3.4,1.1);
\node[font=\scriptsize, poppurple, above] at (3.0,1.12) {$I$};
\draw[thick, popdark] (-1.5,1.1) -- (-2.4,1.1) -- (-2.4,-0.3) -- (-1.5,-0.3);
\node[font=\small\bfseries, align=center] at (0.5,-1.3) {Photon energy $\rightarrow$ electrical energy};
\end{tikzpicture}
\legende{Schematic of a photovoltaic cell: photons with $hf \geq E_{\text{gap}}$ free electrons at the p--n junction; the internal field drives them round the external circuit.}
\end{popfigure}

\courssub{Worked Example: Counting Photons in a Solar Cell}

\begin{exoresolu}[Photons absorbed by a small panel]
A small solar panel of area $\SI{50}{cm^2}$ receives monochromatic light of wavelength $\lambda = \SI{550}{nm}$ with an intensity (power per unit area) of $\SI{200}{W.m^{-2}}$. Calculate (a) the energy of one photon, (b) the number of photons striking the panel each second. Take $h = \SI{6.63e-34}{J.s}$, $c = \SI{3.00e8}{m/s}$.
\tcblower
\textbf{(a) Energy of one photon.}
\[ E = hf = \frac{hc}{\lambda} = \frac{(\SI{6.63e-34}{J.s})(\SI{3.00e8}{m/s})}{\SI{550e-9}{m}} = \SI{3.62e-19}{J} \approx \SI{2.26}{eV}. \]
\textbf{(b) Photon rate.} The power falling on the panel is
\[ P = I \times A = \SI{200}{W.m^{-2}} \times \SI{50e-4}{m^2} = \SI{1.00}{W}. \]
Each photon carries $\SI{3.62e-19}{J}$, so the number arriving per second is
\[ N = \frac{P}{E} = \frac{\SI{1.00}{W}}{\SI{3.62e-19}{J}} = \SI{2.76e18}{s^{-1}}. \]
Even a modest panel processes an enormous photon flux --- which is why the \emph{granular} nature of light is invisible in everyday life.
\end{exoresolu}

\courssub{Synoptic Revision: Wave Model versus Photon Model}

The table below confronts each experimental fact of the photoelectric effect with the predictions of the two models. Learn to reproduce this comparison: it is the backbone of GCE essay answers.

\begin{center}
\begin{tabularx}{\linewidth}{|l|X|X|}
\hline
\rowcolor{popblueL} \textbf{Experimental fact} & \textbf{Wave model predicts} & \textbf{Photon model predicts} \\
\hline
Emission is instantaneous (below $\SI{e-9}{s}$) & Delay while energy accumulates & Immediate: one photon, one electron \\
\hline
Threshold frequency $f_0$ exists & Emission at any frequency if intense enough & No emission if $hf < W_0$ \\
\hline
$E_{k,\max}$ depends only on $f$ (linearly) & $E_{k,\max}$ should grow with intensity & $E_{k,\max} = hf - W_0$ \\
\hline
Number of electrons $\propto$ intensity & (agrees) & More photons $\Rightarrow$ more electrons \\
\hline
$E_{k,\max}$ independent of intensity & Fails & Intensity raises photon \emph{number}, not photon \emph{energy} \\
\hline
\end{tabularx}
\end{center}

\begin{aretenir}[Key points of this section]
\begin{itemize}
\item Interference, diffraction and polarisation prove the \textbf{wave} nature of light; the photoelectric effect, black-body radiation and the Compton effect prove the \textbf{particle} nature.
\item Light is \emph{one} entity with two complementary aspects; the experiment selects the aspect observed.
\item The photon bridges the two: $E = hf$ and $p = h/\lambda$ link particle quantities to wave quantities.
\item Propagation $\rightarrow$ think waves; interaction/detection $\rightarrow$ think photons.
\item Applications: photocells, solar panels, photodiodes, image sensors, photomultipliers.
\end{itemize}
\end{aretenir}

\begin{savaistu}[Exam tip: the duality essay]
GCE Paper 2 essays often ask you to ``discuss the evidence for the dual nature of light''. Structure your answer as follows: (1) define duality; (2) one experiment for the wave aspect (Young's slits, with a sentence on what is observed); (3) one experiment for the particle aspect (photoelectric effect, quoting $hf = W_0 + E_{k,\max}$ and one fact the wave theory cannot explain); (4) conclude that both models are complementary and that $E = hf$ unites them. Four clear paragraphs, four clear marks.
\end{savaistu}

\begin{exercice}[Checking your understanding]
\begin{enumerate}
\item State one phenomenon that demonstrates the wave nature of light and one that demonstrates its particle nature, explaining briefly why each is conclusive.
\item Calculate the momentum of a photon of wavelength $\SI{400}{nm}$.
\item Explain, using the photon model, why a photovoltaic cell does not respond to infrared radiation even when the infrared beam is very intense.
\end{enumerate}
\end{exercice}

\begin{corrige}[Exercise 1]
\textbf{1.} Wave aspect: interference in Young's double-slit experiment --- the fringe pattern (reinforcement and cancellation) requires superposition, which only waves exhibit. Particle aspect: the photoelectric effect --- the existence of a threshold frequency and instantaneous emission show that energy arrives in discrete quanta $hf$, not as a continuous wave.
\textbf{2.} $p = h/\lambda = (\SI{6.63e-34}{J.s})/(\SI{400e-9}{m}) = \SI{1.66e-27}{kg.m.s^{-1}}$.
\textbf{3.} Infrared photons have low frequency, so $hf$ is smaller than the energy gap $E_{\text{gap}}$ of the semiconductor: no electron--hole pair can be created. Raising the intensity only increases the \emph{number} of insufficiently energetic photons --- exactly as in the photoelectric effect below threshold.
\end{corrige}

\newpage

\coursec{Integration activity}

\coursec{Photoelectric Effect: Engineering Application — Automatic Lighting Sensor}

\courssub{Aims of the Activity}
\begin{itemize}
    \item Apply the photon model and Einstein's photoelectric equation to a real-world engineering problem: selecting a photosensitive material for an automatic lighting sensor.
    \item Compute threshold wavelengths from given work functions and interpret them relative to the solar spectrum reaching the ground in Cameroon.
    \item Calculate the maximum kinetic energy of photoelectrons for a specific incident wavelength (green light) and relate it to the sensor's electrical signal (stopping potential).
    \item Construct and interpret an energy-balance diagram illustrating the photoelectric process.
    \item Evaluate the limitations of simple metal photocathodes compared to semiconductor junctions (solar cells) for energy conversion applications.
    \item Develop collaborative problem-solving, data analysis, and technical communication skills.
\end{itemize}

\courssub{Context and Situation}
\textbf{Context:} The administration of \emph{Government Bilingual High School (GBHS) Bamenda} wishes to install an automatic lighting system at the main gate. The system should switch on the security lights at dusk and switch them off at dawn without human intervention. The design team proposes a \emph{photoelectric sensor} (a photocell) mounted on the gate pillar. When sunlight falls on the photocathode, the emitted photoelectrons create a current that keeps a relay open (lights OFF). When the light intensity drops below a threshold (sunset), the photocurrent falls, the relay closes, and the lights turn ON.

As physics consultants for the project, your group is tasked with selecting the most suitable photosensitive metal for the photocathode from a supplier's catalogue. The sensor must operate reliably using natural sunlight in Bamenda (latitude $\approx 6^\circ$ N), where the solar spectrum at ground level ranges approximately from \SI{300}{\nano\meter} (near UV, ozone absorption cut-off) to \SI{2500}{\nano\meter} (infrared), with the visible band (\SI{400}{\nano\meter}--\SI{700}{\nano\meter}) carrying the bulk of the energy. The sensor will be exposed to the elements, so stability is a concern, but the primary selection criterion is \textbf{photosensitivity across the visible spectrum}.

The supplier provides the following data sheet for five candidate metals (work functions $\phi$ given at room temperature):

\begin{center}
\begin{tabularx}{\linewidth}{|l|X|X|}\hline
\rowcolor{popblueL}
\textbf{Metal} & \textbf{Symbol} & \textbf{Work Function $\phi$ (\SI{}{\electronvolt})} \\ \hline
Caesium & Cs & 2.14 \\ \hline
Sodium & Na & 2.28 \\ \hline
Potassium & K & 2.30 \\ \hline
Zinc & Zn & 4.31 \\ \hline
Copper & Cu & 4.65 \\ \hline
\end{tabularx}
\end{center}

\textbf{Useful constants:}
\begin{itemize}
    \item Planck constant $h = \SI{6.63e-34}{\joule\second} = \SI{4.14e-15}{\electronvolt\second}$
    \item Speed of light $c = \SI{3.00e8}{\meter\per\second}$
    \item $hc = \SI{1240}{\electronvolt\nano\meter}$ (useful shortcut for quick calculations)
    \item Elementary charge $e = \SI{1.60e-19}{\coulomb}$
    \item $1\ \mathrm{eV} = \SI{1.60e-19}{\joule}$
\end{itemize}

\courssub{Tasks}
\begin{enumerate}
    \item \textbf{Threshold Wavelengths.} For each of the five metals, calculate the threshold wavelength $\lambda_0$ (in nm) above which no photoelectric emission occurs, regardless of light intensity. Present your results in a table.
    \item \textbf{Spectral Response Analysis.} The solar spectrum at ground level in Bamenda is effectively limited to $\lambda \ge \SI{300}{\nano\meter}$ (ozone absorption) and the useful visible range is \SI{400}{\nano\meter} to \SI{700}{\nano\meter}. Classify each metal as:
    \begin{itemize}
        \item \textbf{Type V:} Responds to a significant portion of the \textbf{visible} spectrum (threshold $\lambda_0 > \SI{400}{\nano\meter}$).
        \item \textbf{Type UV:} Requires \textbf{ultraviolet} light ($\lambda_0 < \SI{400}{\nano\meter}$) for emission; insensitive to visible light.
    \end{itemize}
    Justify your classification using the calculated $\lambda_0$ values.
    \item \textbf{Green Light Response.} The design team tests the sensor with a standard green LED ($\lambda = \SI{550}{\nano\meter}$) to simulate bright daylight. For the metal(s) classified as Type V, calculate the maximum kinetic energy $K_{\text{max}}$ of the emitted photoelectrons (in eV and in joules). Comment on whether the photocurrent generated by green light alone would be sufficient to drive the relay circuit (assume a typical relay requires a stopping potential equivalent to at least \SI{0.10}{\volt} to operate reliably).
    \item \textbf{Material Selection \& Justification.} Based on your calculations and the following practical considerations, recommend \textbf{one} metal for the photocathode:
    \begin{itemize}
        \item Sensitivity to visible light (especially green/yellow peak of solar spectrum).
        \item Magnitude of $K_{\text{max}}$ for green light (signal robustness).
        \item Chemical stability: Alkali metals (Cs, Na, K) oxidise rapidly in air; they must be sealed in an evacuated or inert-gas tube. Transition metals (Zn, Cu) are stable in air but require UV.
    \end{itemize}
    Write a short justification (approx. 100 words).
    \item \textbf{Energy-Balance Diagram.} Draw a labelled energy-level diagram for your chosen metal illuminated by the green LED ($\lambda = \SI{550}{\nano\meter}$). Show: the Fermi level / conduction band, the vacuum level, the work function $\phi$, the incident photon energy $h\nu$, and the maximum kinetic energy $K_{\text{max}}$ of the emitted electron. Indicate the energy balance equation on the diagram.
    \item \textbf{Class Discussion Prompt (Extension).} \emph{``Why do modern solar panels (photovoltaic cells) use semiconductor $p$--$n$ junctions (like silicon) instead of simple metal photocathodes based on the photoelectric effect?''} List three distinct physical or engineering reasons.
\end{enumerate}

\courssub{Model Solutions}

\textbf{Task 1: Threshold Wavelengths}
The threshold wavelength $\lambda_0$ is defined by the condition that the photon energy exactly equals the work function: $h\nu_0 = \phi$, i.e. $\dfrac{hc}{\lambda_0} = \phi$. Using the shortcut $hc = \SI{1240}{\electronvolt\nano\meter}$:
\[
\lambda_0 = \frac{hc}{\phi} = \frac{\SI{1240}{\electronvolt\nano\meter}}{\phi\ (\mathrm{eV})}
\]

\begin{center}
\begin{tabularx}{\linewidth}{|l|X|X|X|}\hline
\rowcolor{popblueL}
\textbf{Metal} & \textbf{$\phi$ (\SI{}{\electronvolt})} & \textbf{Calculation} & \textbf{$\lambda_0$ (\SI{}{\nano\meter})} \\ \hline
Caesium (Cs) & 2.14 & $1240 / 2.14$ & 579 \\ \hline
Sodium (Na) & 2.28 & $1240 / 2.28$ & 544 \\ \hline
Potassium (K) & 2.30 & $1240 / 2.30$ & 539 \\ \hline
Zinc (Zn) & 4.31 & $1240 / 4.31$ & 288 \\ \hline
Copper (Cu) & 4.65 & $1240 / 4.65$ & 267 \\ \hline
\end{tabularx}
\end{center}
\textit{Note: Values are given to three significant figures, consistent with the data provided.}

\textbf{Task 2: Spectral Response Analysis}
\begin{itemize}
    \item \textbf{Type V (Visible response):} \textbf{Caesium ($\lambda_0 = \SI{579}{\nano\meter}$)}, \textbf{Sodium ($\lambda_0 = \SI{544}{\nano\meter}$)}, \textbf{Potassium ($\lambda_0 = \SI{539}{\nano\meter}$)}.
    \emph{Justification:} Their threshold wavelengths lie within or above the visible range ($> \SI{400}{\nano\meter}$). Photons of blue, green, yellow, and (for Cs) orange/red light have sufficient energy ($\lambda < \lambda_0$) to eject electrons. Caesium covers the widest portion of the visible spectrum.
    \item \textbf{Type UV (UV only):} \textbf{Zinc ($\lambda_0 = \SI{288}{\nano\meter}$)}, \textbf{Copper ($\lambda_0 = \SI{267}{\nano\meter}$)}.
    \emph{Justification:} Their threshold wavelengths are below \SI{300}{\nano\meter} (deep UV). Since the atmosphere blocks $\lambda < \SI{300}{\nano\meter}$, \emph{no photoelectrons are emitted by natural sunlight} on these metals. They are unsuitable for an outdoor sunlight sensor.
\end{itemize}

\textbf{Task 3: Green Light Response ($\lambda = \SI{550}{\nano\meter}$)}
Photon energy for green light:
\[
E_{\text{ph}} = \frac{hc}{\lambda} = \frac{\SI{1240}{\electronvolt\nano\meter}}{\SI{550}{\nano\meter}} = \SI{2.255}{\electronvolt}
\]
Einstein's photoelectric equation: $K_{\text{max}} = h\nu - \phi = E_{\text{ph}} - \phi$. We compute only for Type V metals where $E_{\text{ph}} > \phi$.

\begin{center}
\begin{tabularx}{\linewidth}{|l|X|X|X|}\hline
\rowcolor{popblueL}
\textbf{Metal} & \textbf{$\phi$ (\SI{}{\electronvolt})} & \textbf{$K_{\text{max}} = 2.255 - \phi$ (\SI{}{\electronvolt})} & \textbf{$K_{\text{max}}$ (\SI{}{\joule})} \\ \hline
Caesium (Cs) & 2.14 & $2.255 - 2.14 = \mathbf{0.115}$ & $0.115 \times \SI{1.60e-19}{} = \mathbf{\SI{1.84e-20}{}}$ \\ \hline
Sodium (Na) & 2.28 & $2.255 - 2.28 = -0.025$ & \emph{No emission} \\ \hline
Potassium (K) & 2.30 & $2.255 - 2.30 = -0.045$ & \emph{No emission} \\ \hline
\end{tabularx}
\end{center}

\textbf{Analysis:} Only \textbf{Caesium} emits photoelectrons under green light ($\lambda = \SI{550}{\nano\meter}$). Sodium and Potassium have work functions slightly higher than the green photon energy; their thresholds (\SI{544}{\nano\meter} and \SI{539}{\nano\meter}) are shorter than \SI{550}{\nano\meter}, so green photons lack the energy to overcome $\phi$.

For Caesium: $K_{\text{max}} = \SI{0.115}{\electronvolt}$. The corresponding stopping potential is $V_{\text{stop}} = K_{\text{max}}/e = \SI{0.115}{\volt}$.
Since $\SI{0.115}{\volt} > \SI{0.10}{\volt}$ (relay requirement), the photocurrent generated by green light on Caesium \textbf{is sufficient} to keep the relay open (lights OFF) during daylight. The margin is small (\SI{0.015}{\volt}), implying the sensor circuit must be carefully designed (low noise amplifier, shielding from stray capacitance).

\textbf{Task 4: Material Selection \& Justification}
\textbf{Recommended Metal: Caesium (Cs).}

\textbf{Justification:} Caesium is the \emph{only} candidate that responds to green light ($\lambda = \SI{550}{\nano\meter}$), which lies near the peak of the solar spectral irradiance. Its threshold wavelength (\SI{579}{\nano\meter}) covers the entire visible spectrum down to yellow/orange, ensuring maximum utilisation of available sunlight. It yields a positive $K_{\text{max}} = \SI{0.115}{\electronvolt}$ for green light, exceeding the relay's \SI{0.10}{\volt} requirement. Although caesium is highly reactive (oxidises instantly in air), this is standardly managed by sealing the photocathode in an evacuated glass envelope (phototube), a mature and reliable technology for light sensors. Sodium and Potassium fail the green-light test; Zinc and Copper are blind to visible sunlight.

\textbf{Task 5: Energy-Balance Diagram}
\begin{popfigure}
\begin{tikzpicture}[scale=1.1, >=stealth]
    % Styles
    \tikzset{
        level/.style={thick, popdark},
        arrow/.style={->, thick, poporange},
        photon/.style={decorate, decoration={snake, amplitude=1.5pt, segment length=4pt}, popblue, thick, ->},
        label/.style={font=\small, popdark},
        region/.style={fill=popgreenL, opacity=0.2}
    }
    % Coordinates
    \def\yFermi{0}
    \def\yVac{2.14} % 1 cm = 1 eV scale
    \def\yPhotonTop{3.5}
    \def\xMetalLeft{-1.5}
    \def\xMetalRight{2.5}
    \def\xVacRight{5.5}
    \def\xCenter{0.5}
    
    % Metal Region (Conduction Band / Fermi Sea)
    \fill[region] (\xMetalLeft, \yFermi-1.0) rectangle (\xMetalRight, \yFermi);
    \draw[level] (\xMetalLeft, \yFermi) -- (\xMetalRight, \yFermi) node[midway, above=2pt, label] {$E_F$ (Fermi Level)};
    \node[label, align=center] at (\xMetalLeft+0.5, \yFermi-0.5) {Metal \\ (Conduction electrons)};
    
    % Vacuum Level
    \draw[level, dashed] (\xMetalRight, \yVac) -- (\xVacRight, \yVac) node[midway, above=2pt, label] {$E_{\text{vac}}$ (Vacuum Level)};
    \node[label, align=center] at (\xVacRight-0.5, \yVac+0.4) {Free electron \\ (at rest in vacuum)};
    
    % Work Function arrow
    \draw[arrow] (\xCenter, \yFermi) -- (\xCenter, \yVac) node[midway, right=4pt, label] {$\phi = \SI{2.14}{\electronvolt}$};
    
    % Incident Photon (absorbed by electron at Fermi level)
    \draw[photon] (\xCenter, \yPhotonTop) -- (\xCenter, \yFermi) node[midway, left=4pt, label] {$h\nu = \SI{2.255}{\electronvolt}$};
    \node[label, popblue, align=center] at (\xCenter, \yPhotonTop+0.4) {Green photon\\ $\lambda=\SI{550}{\nano\meter}$};
    
    % Energy gained by electron (virtual line)
    \draw[popmuted, dashed, thin] (\xCenter, \yFermi) -- (\xCenter, \yPhotonTop);
    \node[label, popmuted, align=left] at (\xCenter+0.8, \yPhotonTop-0.2) {$E_F + h\nu$};
    
    % Emitted Electron with K_max
    \draw[arrow, poporange] (\xCenter, \yVac) -- (\xVacRight-0.5, \yVac+0.6) node[midway, above, label] {$K_{\text{max}} = \SI{0.115}{\electronvolt}$};
    \node[label, poporange, align=center] at (\xVacRight-0.5, \yVac+1.0) {Photoelectron \\ emitted};
    
    % Energy Balance Equation Box
    \node[draw, popblue, thick, rounded corners, fill=popblueL, align=left, label, anchor=west] at (\xVacRight+0.3, \yFermi) {
        \textbf{Einstein's Equation:}\\
        $h\nu = \phi + K_{\text{max}}$\\
        $\SI{2.255}{\electronvolt} = \SI{2.14}{\electronvolt} + \SI{0.115}{\electronvolt}$
    };
    
    % Energy Axis
    \draw[<->, thin, popmuted] (\xMetalLeft-0.8, \yFermi-1.0) -- (\xMetalLeft-0.8, \yPhotonTop+0.2) node[midway, left=6pt, label, rotate=90] {Energy / eV};
    % Scale ticks
    \foreach \y/\lab in {0/0, 2.14/2.14, 3.5/3.5} {
        \draw[thin, popmuted] (\xMetalLeft-0.9, \y) -- (\xMetalLeft-0.7, \y);
        \node[label, popmuted, anchor=east] at (\xMetalLeft-0.95, \y) {\lab};
    }
\end{tikzpicture}
\legende{Energy-balance diagram for Caesium photocathode illuminated by green light ($\lambda = \SI{550}{\nano\meter}$). An electron at the Fermi level $E_F$ absorbs a photon of energy $h\nu$. If $h\nu > \phi$, the electron escapes the metal; its maximum kinetic energy in vacuum is $K_{\text{max}} = h\nu - \phi$.}
\end{popfigure}

\textbf{Task 6: Class Discussion — Why Semiconductors (Solar Cells) instead of Metal Photocathodes?}
Three key reasons:
\begin{enumerate}
    \item \textbf{Quantum Efficiency \& Internal Photoelectric Effect:} In a metal, the photoelectric effect is a \emph{surface} phenomenon: only electrons within $\sim \SI{1}{\nano\meter}$ of the surface can escape without losing energy to collisions. Most absorbed photons simply heat the metal (low quantum efficiency, typically $< 0.1\%$). In a semiconductor $p$--$n$ junction, the \emph{internal} photoelectric effect creates electron-hole pairs throughout the bulk depletion region (microns thick). The built-in electric field separates them before recombination, achieving quantum efficiencies $> 80\%$.
    \item \textbf{Voltage \& Power Output:} A metal photocathode emits electrons into vacuum; the external circuit must collect them against a space-charge limit, yielding very low current densities (microamps per $\mathrm{cm}^2$) and low voltage ($V_{\text{stop}} \sim \phi/e$). A $p$--$n$ junction generates a photovoltage ($\sim 0.6\text{--}0.7\ \mathrm{V}$ for Si) and delivers useful power (mA/$\mathrm{cm}^2$ at $\sim 0.5\ \mathrm{V}$) directly to a load, functioning as a \emph{generator}, not just a switch.
    \item \textbf{Material Stability \& Bandgap Engineering:} Alkali metals (low $\phi$) oxidise instantly, requiring expensive vacuum tubes. Semiconductors (Si, GaAs, perovskites) are solid, stable in air (with passivation), and their bandgap $E_g$ (analogous to $\phi$) can be engineered ($1.1\text{--}1.8\ \mathrm{eV}$) to match the solar spectrum optimally (Shockley-Queisser limit), unlike fixed metal work functions.
\end{enumerate}

\begin{aretenir}[Key Takeaways for the GCE Exam]
\begin{itemize}
\item \textbf{Threshold condition:} $\lambda_0 = hc/\phi$. Use $hc = \SI{1240}{\electronvolt\nano\meter}$ for speed in eV/nm calculations.
\item \textbf{Einstein's Equation:} $K_{\text{max}} = h\nu - \phi = hc/\lambda - \phi$. \textbf{Always check $h\nu > \phi$ first!} If not, $K_{\text{max}} = 0$ (no emission).
\item \textbf{Stopping Potential:} $eV_{\text{stop}} = K_{\text{max}}$. This links the quantum result to a measurable voltage.
\item \textbf{Metal vs Semiconductor:} Metals $\rightarrow$ surface effect, low yield, vacuum tube needed. Semiconductors $\rightarrow$ bulk effect, high yield, solid-state, power generation.
\item \textbf{Units:} Convert eV to Joules ($\times \SI{1.60e-19}{\joule/\electronvolt}$) only when explicitly asked for SI units. Keep energies in eV for photoelectric calculations.
\end{itemize}
\end{aretenir}

\newpage

\coursec{Methods and worked examples}

\courssub{Methods and Worked Examples}

\begin{methode}[Calculating Photon Energy and Relating to Frequency/Wavelength]
\textbf{Step 1:} Identify the given data: frequency $f$ (in Hz) or wavelength $\lambda$ (in m).\par
\textbf{Step 2:} Recall the photon energy formulas: $E = hf = \dfrac{hc}{\lambda}$, where $h = 6.63\times10^{-34}\ \mathrm{J\,s}$ and $c = 3.00\times10^8\ \mathrm{m/s}$.\par
\textbf{Step 3:} If $\lambda$ is given in nanometres (nm), convert to metres: $1\ \mathrm{nm} = 10^{-9}\ \mathrm{m}$.\par
\textbf{Step 4:} Substitute values and compute $E$ in joules. Often the answer is required in electronvolts (eV): $1\ \mathrm{eV} = 1.60\times10^{-19}\ \mathrm{J}$.\par
\textbf{Step 5:} Check significant figures (usually 3 s.f. for GCE) and include the correct unit.
\end{methode}

\begin{exoresolu}[Photon Energy of Visible Light]
A helium-neon laser emits red light of wavelength $\lambda = 632.8\ \mathrm{nm}$. Calculate the energy of a single photon in joules and in electronvolts.
\tcblower
\textbf{Solution:}\par
\underline{Given:} $\lambda = 632.8\ \mathrm{nm} = 632.8\times10^{-9}\ \mathrm{m} = 6.328\times10^{-7}\ \mathrm{m}$.\par
\underline{Constants:} $h = 6.63\times10^{-34}\ \mathrm{J\,s}$, $c = 3.00\times10^8\ \mathrm{m/s}$, $1\ \mathrm{eV} = 1.60\times10^{-19}\ \mathrm{J}$.\par
\underline{Formula:} $E = \dfrac{hc}{\lambda}$.\par
\[
E = \frac{(6.63\times10^{-34})(3.00\times10^8)}{6.328\times10^{-7}}
   = \frac{1.989\times10^{-25}}{6.328\times10^{-7}}
   = 3.14\times10^{-19}\ \mathrm{J}.
\]
Convert to eV:
\[
E = \frac{3.14\times10^{-19}\ \mathrm{J}}{1.60\times10^{-19}\ \mathrm{J/eV}} = 1.96\ \mathrm{eV}.
\]
\underline{Answer:} $E = 3.14\times10^{-19}\ \mathrm{J}$ or $1.96\ \mathrm{eV}$.
\end{exoresolu}

\begin{methode}[Applying Einstein's Photoelectric Equation]
\textbf{Step 1:} Write down Einstein's equation: $hf = \phi + K_{\max}$, where $\phi = hf_0$ is the work function and $K_{\max} = eV_s$ is the maximum kinetic energy of emitted electrons ($V_s$ = stopping potential).\par
\textbf{Step 2:} Identify knowns: frequency $f$ (or wavelength $\lambda$), work function $\phi$ (or threshold frequency $f_0$), stopping potential $V_s$.\par
\textbf{Step 3:} Convert all energies to the same unit (usually joules or eV). If $\phi$ is given in eV, convert to joules: $\phi(\mathrm{J}) = \phi(\mathrm{eV}) \times 1.60\times10^{-19}$.\par
\textbf{Step 4:} Rearrange the equation to solve for the unknown: $K_{\max} = hf - \phi$, $V_s = \dfrac{hf - \phi}{e}$, $f_0 = \dfrac{\phi}{h}$, etc.\par
\textbf{Step 5:} Substitute numerical values, compute, and express the answer with correct units (J, eV, V, Hz).
\end{methode}

\begin{exoresolu}[Stopping Potential and Threshold Frequency]
The work function of sodium is $\phi = 2.28\ \mathrm{eV}$. Light of wavelength $400\ \mathrm{nm}$ is incident on a sodium surface. Calculate: (a) the maximum kinetic energy of emitted photoelectrons in eV, (b) the stopping potential, (c) the threshold frequency for sodium.
\tcblower
\textbf{Solution:}\par
\underline{Given:} $\phi = 2.28\ \mathrm{eV}$, $\lambda = 400\ \mathrm{nm} = 4.00\times10^{-7}\ \mathrm{m}$.\par
\underline{Constants:} $h = 6.63\times10^{-34}\ \mathrm{J\,s}$, $c = 3.00\times10^8\ \mathrm{m/s}$, $e = 1.60\times10^{-19}\ \mathrm{C}$.\par
\underline{(a) Photon energy:}
\[
E = \frac{hc}{\lambda} = \frac{(6.63\times10^{-34})(3.00\times10^8)}{4.00\times10^{-7}} = 4.97\times10^{-19}\ \mathrm{J}.
\]
Convert to eV: $E = \dfrac{4.97\times10^{-19}}{1.60\times10^{-19}} = 3.11\ \mathrm{eV}$.\par
Maximum kinetic energy:
\[
K_{\max} = E - \phi = 3.11\ \mathrm{eV} - 2.28\ \mathrm{eV} = 0.83\ \mathrm{eV}.
\]
\underline{(b) Stopping potential:} $K_{\max} = eV_s \Rightarrow V_s = \dfrac{K_{\max}}{e} = 0.83\ \mathrm{V}$ (a kinetic energy of $1\ \mathrm{eV}$ corresponds to a stopping potential of $1\ \mathrm{V}$ for an electron).\par
\underline{(c) Threshold frequency:} $\phi = hf_0 \Rightarrow f_0 = \dfrac{\phi}{h}$.
Convert $\phi$ to joules: $\phi = 2.28 \times 1.60\times10^{-19} = 3.65\times10^{-19}\ \mathrm{J}$.
\[
f_0 = \frac{3.65\times10^{-19}}{6.63\times10^{-34}} = 5.50\times10^{14}\ \mathrm{Hz}.
\]
\underline{Answers:} (a) $0.83\ \mathrm{eV}$, (b) $0.83\ \mathrm{V}$, (c) $5.50\times10^{14}\ \mathrm{Hz}$.
\end{exoresolu}

\begin{attention}[Units Trap in Photoelectric Calculations]
The most frequent GCE error is mixing joules and electronvolts. If you use $h = 6.63\times10^{-34}\ \mathrm{J\,s}$, then $hf$ is in \emph{joules}: the work function must also be in joules before subtracting. Alternatively, convert the photon energy to eV first and work entirely in eV. Never subtract an energy in eV from an energy in J.
\end{attention}

\begin{methode}[Graphical Determination of Planck's Constant and Work Function]
\textbf{Step 1:} Recall the linear form of Einstein's equation: $eV_s = hf - \phi$ or $V_s = \dfrac{h}{e}f - \dfrac{\phi}{e}$.\par
\textbf{Step 2:} Plot $V_s$ (y-axis) against $f$ (x-axis). The graph is a straight line with slope $m = \dfrac{h}{e}$ and y-intercept $c = -\dfrac{\phi}{e}$.\par
\textbf{Step 3:} Choose two well-separated points on the best-fit line to calculate the slope: $m = \dfrac{\Delta V_s}{\Delta f}$.\par
\textbf{Step 4:} Compute Planck's constant: $h = m \times e$.\par
\textbf{Step 5:} Determine the work function from the intercept: $\phi = -c \times e$ (or from the threshold frequency $f_0$ where $V_s = 0$: $\phi = hf_0$).\par
\textbf{Step 6:} Express $h$ in $\mathrm{J\,s}$ and $\phi$ in $\mathrm{J}$ or $\mathrm{eV}$.
\end{methode}

\begin{exoresolu}[Finding $h$ and $\phi$ from Experimental Data]
In a photoelectric experiment, the following stopping potentials $V_s$ were measured for different frequencies $f$ of incident light:

\begin{center}
\begin{tabular}{|c|c|c|c|c|}
\hline
$f\ (10^{14}\ \mathrm{Hz})$ & 5.0 & 6.0 & 7.0 & 8.0 \\ \hline
$V_s\ (\mathrm{V})$ & 0.25 & 0.66 & 1.07 & 1.48 \\ \hline
\end{tabular}
\end{center}
Plot $V_s$ vs $f$ and determine (a) Planck's constant $h$, (b) the work function $\phi$ in eV.
\tcblower
\textbf{Solution:}\par
\underline{Graph:} Plot the points $(5.0\,;\,0.25)$, $(6.0\,;\,0.66)$, $(7.0\,;\,1.07)$, $(8.0\,;\,1.48)$ and draw the best-fit straight line.\par
\underline{Slope calculation:} Use the points $(5.0\,;\,0.25)$ and $(8.0\,;\,1.48)$.
\[
m = \frac{1.48 - 0.25}{(8.0 - 5.0)\times10^{14}} = \frac{1.23}{3.0\times10^{14}} = 4.10\times10^{-15}\ \mathrm{V\,s}.
\]
\underline{(a) Planck's constant:} $h = m \times e = (4.10\times10^{-15}) \times (1.60\times10^{-19}) = 6.56\times10^{-34}\ \mathrm{J\,s}$.\par
\underline{(b) Work function:} Find the y-intercept $c$. Using the point $(5.0\,;\,0.25)$:
\[
c = V_s - m f = 0.25 - (4.10\times10^{-15})(5.0\times10^{14}) = 0.25 - 2.05 = -1.80\ \mathrm{V}.
\]
Then $\phi = -c \times e = 1.80 \times 1.60\times10^{-19} = 2.88\times10^{-19}\ \mathrm{J}$.
In eV: $\phi = \dfrac{2.88\times10^{-19}}{1.60\times10^{-19}} = 1.80\ \mathrm{eV}$.\par
\underline{Answers:} (a) $h = 6.56\times10^{-34}\ \mathrm{J\,s}$ (close to the accepted value $6.63\times10^{-34}\ \mathrm{J\,s}$), (b) $\phi = 1.80\ \mathrm{eV}$.
\end{exoresolu}

\begin{methode}[Photoelectric Effect with Two Frequencies / Unknown Frequency]
\textbf{Step 1:} Write Einstein's equation for each frequency: $hf_1 = \phi + eV_{s1}$ and $hf_2 = \phi + eV_{s2}$.\par
\textbf{Step 2:} Subtract the two equations to eliminate $\phi$: $h(f_2 - f_1) = e(V_{s2} - V_{s1})$.\par
\textbf{Step 3:} Solve for the unknown frequency or the work function.\par
\textbf{Step 4:} If the threshold frequency $f_0$ is involved, use $\phi = hf_0$.\par
\textbf{Step 5:} Check consistency: the unknown frequency must be greater than $f_0$ for photoelectric emission to occur.
\end{methode}

\begin{exoresolu}[Determining an Unknown Frequency]
A metal surface has a work function $\phi = 2.00\ \mathrm{eV}$. When illuminated by light of frequency $f_1 = 8.00\times10^{14}\ \mathrm{Hz}$, the stopping potential is $V_{s1} = 1.32\ \mathrm{V}$. A second light source of unknown frequency $f_2$ gives a stopping potential $V_{s2} = 2.50\ \mathrm{V}$. Find $f_2$.
\tcblower
\textbf{Solution:}\par
\underline{Method 1: Direct use of Einstein's equation for the second source.}\par
Convert $\phi$ to joules: $\phi = 2.00 \times 1.60\times10^{-19} = 3.20\times10^{-19}\ \mathrm{J}$.\par
Consistency check with the first source:
\[
hf_1 = (6.63\times10^{-34})(8.00\times10^{14}) = 5.30\times10^{-19}\ \mathrm{J},
\]
\[
\phi + eV_{s1} = 3.20\times10^{-19} + (1.60\times10^{-19})(1.32) = 3.20\times10^{-19} + 2.11\times10^{-19} = 5.31\times10^{-19}\ \mathrm{J}.
\]
The data are consistent (small rounding difference).\par
For the second source: $hf_2 = \phi + eV_{s2}$.
\[
hf_2 = 3.20\times10^{-19} + (1.60\times10^{-19})(2.50) = 3.20\times10^{-19} + 4.00\times10^{-19} = 7.20\times10^{-19}\ \mathrm{J}.
\]
\[
f_2 = \frac{7.20\times10^{-19}}{6.63\times10^{-34}} = 1.09\times10^{15}\ \mathrm{Hz}.
\]
\underline{Method 2: Subtraction (eliminates $\phi$).}
\[
h(f_2 - f_1) = e(V_{s2} - V_{s1}) \Rightarrow f_2 = f_1 + \frac{e(V_{s2} - V_{s1})}{h}.
\]
\[
f_2 = 8.00\times10^{14} + \frac{(1.60\times10^{-19})(2.50 - 1.32)}{6.63\times10^{-34}}
     = 8.00\times10^{14} + \frac{1.89\times10^{-19}}{6.63\times10^{-34}}
     = 8.00\times10^{14} + 2.85\times10^{14}
     = 1.08\times10^{15}\ \mathrm{Hz}.
\]
Both methods agree (the tiny difference comes from rounding).\par
\underline{Answer:} $f_2 \approx 1.09\times10^{15}\ \mathrm{Hz}$.
\end{exoresolu}

\begin{methode}[Wave-Particle Duality: de Broglie Wavelength]
\textbf{Step 1:} Recall de Broglie's hypothesis: a particle of momentum $p$ has an associated wavelength $\lambda = \dfrac{h}{p}$.\par
\textbf{Step 2:} For an electron accelerated from rest through a potential difference $V$, its kinetic energy is $K = eV$. If non-relativistic ($V \ll 10^5\ \mathrm{V}$), $p = \sqrt{2m_e K} = \sqrt{2m_e eV}$.\par
\textbf{Step 3:} Substitute into the de Broglie formula: $\lambda = \dfrac{h}{\sqrt{2m_e eV}}$.\par
\textbf{Step 4:} Compute $\lambda$ in metres; it is often expressed in nanometres (nm) or picometres (pm).
\end{methode}

\begin{exoresolu}[de Broglie Wavelength of an Electron]
Calculate the de Broglie wavelength of an electron accelerated from rest through a potential difference of $150\ \mathrm{V}$. (Take $m_e = 9.11\times10^{-31}\ \mathrm{kg}$, $e = 1.60\times10^{-19}\ \mathrm{C}$, $h = 6.63\times10^{-34}\ \mathrm{J\,s}$.)
\tcblower
\textbf{Solution:}\par
\underline{Kinetic energy:} $K = eV = (1.60\times10^{-19})(150) = 2.40\times10^{-17}\ \mathrm{J}$.\par
\underline{Momentum:}
\[
p = \sqrt{2m_e K} = \sqrt{2(9.11\times10^{-31})(2.40\times10^{-17})} = \sqrt{4.37\times10^{-47}} = 6.61\times10^{-24}\ \mathrm{kg\,m\,s^{-1}}.
\]
\underline{de Broglie wavelength:}
\[
\lambda = \frac{h}{p} = \frac{6.63\times10^{-34}}{6.61\times10^{-24}} = 1.00\times10^{-10}\ \mathrm{m} = 0.100\ \mathrm{nm} = 100\ \mathrm{pm}.
\]
\underline{Useful shortcut:} $\lambda = \dfrac{h}{\sqrt{2m_e eV}} = \dfrac{1.23}{\sqrt{V}}\ \mathrm{nm}$ (with $V$ in volts). For $V = 150\ \mathrm{V}$: $\lambda = \dfrac{1.23}{\sqrt{150}} = 0.100\ \mathrm{nm}$.\par
\underline{Answer:} $\lambda = 1.00\times10^{-10}\ \mathrm{m}$ (or $0.100\ \mathrm{nm}$).
\end{exoresolu}

\begin{methode}[Intensity, Saturation Current and Stopping Potential]
\textbf{Step 1:} Effect of intensity: increasing the light intensity increases the number of photons per second, hence the \emph{saturation photocurrent} increases proportionally, but the \emph{stopping potential} $V_s$ (and thus $K_{\max}$) remains unchanged because the photon energy $hf$ is unchanged.\par
\textbf{Step 2:} Effect of frequency: increasing $f$ increases the photon energy, thus increases $K_{\max}$ and $V_s$; if the intensity is kept constant, the number of photons per second decreases, so the saturation current decreases.\par
\textbf{Step 3:} Saturation current $I_{\text{sat}} = n e$, where $n$ is the number of photoelectrons collected per second.\par
\textbf{Step 4:} For a source of power $P$ emitting monochromatic light, the number of photons per second is $N = \dfrac{P}{hf}$. If the quantum efficiency is $\eta$ (fraction of photons that eject an electron), then $I_{\text{sat}} = \eta e N = \eta e \dfrac{P}{hf}$.\par
\textbf{Step 5:} Apply these ideas to predict the effect of changes in intensity or frequency on $I_{\text{sat}}$ and $V_s$.
\end{methode}

\begin{exoresolu}[Intensity and Frequency Effects on Photocurrent]
A monochromatic light source of power $P = 2.0\ \mathrm{mW}$ and wavelength $\lambda = 500\ \mathrm{nm}$ illuminates a photocathode of work function $\phi = 2.0\ \mathrm{eV}$. The quantum efficiency is $0.10\%$ (one photoelectron per 1000 incident photons). Calculate: (a) the saturation photocurrent, (b) the stopping potential. If the intensity is doubled while keeping the same wavelength, state the new saturation current and stopping potential.
\tcblower
\textbf{Solution:}\par
\underline{Constants:} $h = 6.63\times10^{-34}\ \mathrm{J\,s}$, $c = 3.00\times10^8\ \mathrm{m/s}$, $e = 1.60\times10^{-19}\ \mathrm{C}$.\par
\underline{Photon energy:} $E = \dfrac{hc}{\lambda} = \dfrac{(6.63\times10^{-34})(3.00\times10^8)}{500\times10^{-9}} = 3.98\times10^{-19}\ \mathrm{J} = 2.49\ \mathrm{eV}$.\par
\underline{(a) Number of incident photons per second:} $N = \dfrac{P}{E} = \dfrac{2.0\times10^{-3}}{3.98\times10^{-19}} = 5.03\times10^{15}\ \mathrm{s^{-1}}$.\par
\underline{Photoelectrons per second:} $n = \eta N = 0.0010 \times 5.03\times10^{15} = 5.03\times10^{12}\ \mathrm{s^{-1}}$.\par
\underline{Saturation current:} $I_{\text{sat}} = n e = (5.03\times10^{12})(1.60\times10^{-19}) = 8.0\times10^{-7}\ \mathrm{A} = 0.80\ \mu\mathrm{A}$.\par
\underline{(b) Stopping potential:} $K_{\max} = E - \phi = 2.49\ \mathrm{eV} - 2.0\ \mathrm{eV} = 0.49\ \mathrm{eV}$, so $V_s = 0.49\ \mathrm{V}$.\par
\underline{Effect of doubling the intensity:} the power becomes $4.0\ \mathrm{mW}$. The photon energy is unchanged, so $V_s$ stays at $0.49\ \mathrm{V}$. The number of photons per second doubles, so $n$ doubles and $I_{\text{sat}}$ doubles to $1.6\ \mu\mathrm{A}$.\par
\underline{Answers:} (a) $0.80\ \mu\mathrm{A}$, (b) $0.49\ \mathrm{V}$. After doubling the intensity: $I_{\text{sat}} = 1.6\ \mu\mathrm{A}$, $V_s = 0.49\ \mathrm{V}$ (unchanged).
\end{exoresolu}

\begin{methode}[Photoelectric Effect in Practical Devices: Photocells and Solar Cells]
\textbf{Step 1:} Identify the type of device: vacuum photocell (photoemissive) or semiconductor photodiode/solar cell (photovoltaic).\par
\textbf{Step 2:} For a vacuum photocell: current flows only if $f > f_0$; the saturation current is proportional to the light intensity; the stopping potential measures $K_{\max}$.\par
\textbf{Step 3:} For a solar cell (p-n junction): photons create electron-hole pairs; the built-in electric field of the junction separates them, generating a voltage and a current without any external stopping potential.\par
\textbf{Step 4:} In problems, you may be asked for the number of photons needed for a given current, using $I = \eta e \dfrac{P_{\text{light}}}{hf}$, or for an efficiency $\eta = \dfrac{P_{\text{out}}}{P_{\text{in}}}$.
\end{methode}

\begin{exoresolu}[Solar Cell Current Calculation]
A solar cell of area $10\ \mathrm{cm}^2$ is illuminated by sunlight of intensity $800\ \mathrm{W\,m^{-2}}$. The cell has a quantum efficiency of $20\%$ for photons of wavelength $600\ \mathrm{nm}$ (assume, for simplicity, that all the sunlight is at this wavelength). Calculate the short-circuit current $I_{sc}$ if every absorbed photon generates one collected electron-hole pair.
\tcblower
\textbf{Solution:}\par
\underline{Given:} $A = 10\ \mathrm{cm}^2 = 1.0\times10^{-3}\ \mathrm{m}^2$; intensity $I = 800\ \mathrm{W\,m^{-2}}$; $\lambda = 600\ \mathrm{nm} = 6.00\times10^{-7}\ \mathrm{m}$; $\eta = 0.20$.\par
\underline{Power incident on the cell:} $P_{\text{in}} = I \times A = 800 \times 1.0\times10^{-3} = 0.80\ \mathrm{W}$.\par
\underline{Photon energy:} $E_{\text{ph}} = \dfrac{hc}{\lambda} = \dfrac{(6.63\times10^{-34})(3.00\times10^8)}{6.00\times10^{-7}} = 3.32\times10^{-19}\ \mathrm{J}$.\par
\underline{Incident photons per second:} $N_{\text{in}} = \dfrac{P_{\text{in}}}{E_{\text{ph}}} = \dfrac{0.80}{3.32\times10^{-19}} = 2.41\times10^{18}\ \mathrm{s^{-1}}$.\par
\underline{Collected electrons per second:} $n = \eta N_{\text{in}} = 0.20 \times 2.41\times10^{18} = 4.83\times10^{17}\ \mathrm{s^{-1}}$.\par
\underline{Short-circuit current:} $I_{sc} = n e = (4.83\times10^{17})(1.60\times10^{-19}) = 7.7\times10^{-2}\ \mathrm{A} = 77\ \mathrm{mA}$.\par
\underline{Answer:} $I_{sc} \approx 77\ \mathrm{mA}$.
\end{exoresolu}

\begin{methode}[Comprehensive Problem: Graph Analysis and de Broglie Wavelength]
\textbf{Step 1:} Read the problem carefully; identify all given data and what is asked.\par
\textbf{Step 2:} Use Einstein's equation $hf = \phi + K_{\max}$ with $K_{\max} = eV_s$.\par
\textbf{Step 3:} If a graph of $V_s$ against $f$ is provided or to be drawn, determine the slope and intercept to find $h$ and $\phi$.\par
\textbf{Step 4:} For the emitted electrons, if the de Broglie wavelength is asked, compute $K_{\max}$ first, then $p = \sqrt{2m_e K_{\max}}$, then $\lambda = h/p$.\par
\textbf{Step 5:} Check units at every step; convert eV to J when using $h$ in $\mathrm{J\,s}$.\par
\textbf{Step 6:} Present final answers with appropriate significant figures and units.
\end{methode}

\begin{exoresolu}[Complete Analysis: From Graph to de Broglie Wavelength]
In a photoelectric experiment on a metal, the following data are obtained:

\begin{center}
\begin{tabular}{|c|c|c|c|c|c|}
\hline
Frequency $f\ (10^{14}\ \mathrm{Hz})$ & 4.0 & 5.0 & 6.0 & 7.0 & 8.0 \\ \hline
Stopping potential $V_s\ (\mathrm{V})$ & 0.10 & 0.52 & 0.94 & 1.36 & 1.78 \\ \hline
\end{tabular}
\end{center}
(a) Plot $V_s$ against $f$ and determine Planck's constant $h$ and the work function $\phi$ (in eV).\par
(b) For the frequency $f = 7.0\times10^{14}\ \mathrm{Hz}$, calculate the de Broglie wavelength of the fastest emitted photoelectrons.
\tcblower
\textbf{Solution:}\par
\underline{(a) Graph and slope:} plot the points; the best-fit line passes approximately through $(4.0\,;\,0.10)$ and $(8.0\,;\,1.78)$.\par
Slope: $m = \dfrac{1.78 - 0.10}{(8.0 - 4.0)\times10^{14}} = \dfrac{1.68}{4.0\times10^{14}} = 4.20\times10^{-15}\ \mathrm{V\,s}$.\par
$h = m e = (4.20\times10^{-15})(1.60\times10^{-19}) = 6.72\times10^{-34}\ \mathrm{J\,s}$.\par
Intercept: using the point $(4.0\,;\,0.10)$:
\[
c = 0.10 - (4.20\times10^{-15})(4.0\times10^{14}) = 0.10 - 1.68 = -1.58\ \mathrm{V}.
\]
$\phi = -c\, e = 1.58 \times 1.60\times10^{-19} = 2.53\times10^{-19}\ \mathrm{J} = \dfrac{2.53\times10^{-19}}{1.60\times10^{-19}} = 1.58\ \mathrm{eV}$.\par
\underline{(b) de Broglie wavelength for $f = 7.0\times10^{14}\ \mathrm{Hz}$:}\par
Photon energy: $E = hf = (6.72\times10^{-34})(7.0\times10^{14}) = 4.70\times10^{-19}\ \mathrm{J}$.\par
$K_{\max} = E - \phi = 4.70\times10^{-19} - 2.53\times10^{-19} = 2.17\times10^{-19}\ \mathrm{J}$.\par
Momentum: $p = \sqrt{2m_e K_{\max}} = \sqrt{2(9.11\times10^{-31})(2.17\times10^{-19})} = \sqrt{3.95\times10^{-49}} = 6.29\times10^{-25}\ \mathrm{kg\,m\,s^{-1}}$.\par
de Broglie wavelength: $\lambda = \dfrac{h}{p} = \dfrac{6.72\times10^{-34}}{6.29\times10^{-25}} = 1.07\times10^{-9}\ \mathrm{m} = 1.07\ \mathrm{nm}$.\par
\underline{Answers:} (a) $h = 6.72\times10^{-34}\ \mathrm{J\,s}$, $\phi = 1.58\ \mathrm{eV}$; (b) $\lambda = 1.07\ \mathrm{nm}$.
\end{exoresolu}

\begin{attention}[Common Mistakes in This Chapter]
\begin{itemize}
\item Using $E = hf$ with $f$ in $10^{14}\ \mathrm{Hz}$ but forgetting the factor $10^{14}$ when substituting.
\item Mixing joules and electronvolts inside Einstein's equation.
\item Believing that increasing the intensity increases $K_{\max}$: only the \emph{number} of photoelectrons increases; $K_{\max}$ depends only on the frequency.
\item Forgetting that no emission occurs if $f < f_0$, however intense the light.
\item In de Broglie calculations, using $K$ in eV inside $p = \sqrt{2m_e K}$: convert to joules first.
\end{itemize}
\end{attention}

\newpage

\coursec{Exercises}

\courssub{I check my knowledge}

\begin{exercice}[MCQ: Photoelectric effect basics]
A monochromatic light of frequency \(f\) is incident on a metal surface of work function \(\phi\). Which of the following statements is correct?
\begin{enumerate}
\item The maximum kinetic energy of emitted electrons depends on the intensity of the incident light.
\item No electrons are emitted if \(f < f_0\), where \(f_0 = \phi/h\), regardless of the light intensity.
\item The stopping potential \(V_s\) is directly proportional to the intensity of the light.
\item The photoelectric current is independent of the frequency of the incident light.
\end{enumerate}
\end{exercice}

\begin{exercice}[True/False with justification]
State whether each of the following statements is true or false. Justify your answer briefly.
\begin{enumerate}
\item Increasing the intensity of incident light increases the maximum kinetic energy of photoelectrons.
\item The photoelectric effect demonstrates the wave nature of light.
\item The threshold frequency of a metal depends on the intensity of the incident light.
\item In Einstein's photoelectric equation, the term \(hf\) represents the energy of a photon.
\end{enumerate}
\end{exercice}

\begin{exercice}[Definitions and key concepts]
Define the following terms:
\begin{enumerate}
\item Work function of a metal.
\item Threshold frequency.
\item Stopping potential.
\item Photon.
\end{enumerate}
\end{exercice}

\begin{exercice}[Direct calculation: Photon energy and photoelectric emission]
Light of wavelength \(400\ \mathrm{nm}\) falls on a sodium surface of work function \(2.28\ \mathrm{eV}\).
\begin{enumerate}
\item Calculate the energy of a single photon in joules and in electronvolts.
\item Determine the maximum kinetic energy of the emitted photoelectrons (in eV).
\item Find the stopping potential required to prevent the most energetic electrons from reaching the anode.
\end{enumerate}
\end{exercice}

\courssub{I practise}

\begin{exercice}[Threshold wavelength and work function]
The threshold wavelength of a certain metal is \(540\ \mathrm{nm}\).
\begin{enumerate}
\item Determine its work function in joules and in electronvolts.
\item Will visible red light of wavelength \(700\ \mathrm{nm}\) liberate electrons from this metal? Explain.
\item If ultraviolet light of wavelength \(300\ \mathrm{nm}\) is used, calculate the maximum kinetic energy of the emitted electrons (in eV).
\end{enumerate}
\end{exercice}

\begin{exercice}[Einstein's equation: Graphical analysis]
In a photoelectric experiment, the stopping potential \(V_s\) is measured for three different frequencies of incident light. The results are:

\begin{center}
\begin{tabular}{|c|c|c|c|}
\hline
Frequency \(f\) (\(\times 10^{14}\ \mathrm{Hz}\)) & 5.0 & 7.0 & 9.0 \\
\hline
Stopping potential \(V_s\) (V) & 0.45 & 1.28 & 2.10 \\
\hline
\end{tabular}
\end{center}

\begin{enumerate}
\item Plot a graph of \(V_s\) against \(f\) (you may describe the expected linear relationship).
\item From the graph, deduce the value of Planck's constant \(h\) (in \(\mathrm{J\,s}\)) and the work function \(\phi\) (in eV) of the metal.
\item Compare your value of \(h\) with the accepted value \(6.63 \times 10^{-34}\ \mathrm{J\,s}\) and comment on possible sources of error.
\end{enumerate}
\end{exercice}

\begin{exercice}[Photon flux and photocurrent]
A \(60\ \mathrm{W}\) lamp emits \(2\%\) of its power as visible light of average wavelength \(550\ \mathrm{nm}\). A photocell of area \(1\ \mathrm{cm}^2\) is placed \(2\ \mathrm{m}\) from the lamp. Assume the light spreads uniformly in all directions.
\begin{enumerate}
\item Calculate the number of photons emitted by the lamp per second in the visible range.
\item Determine the intensity of the visible light at the position of the photocell.
\item Estimate the number of photons striking the photocell each second.
\item If each photon liberates one electron (quantum efficiency = 100\%), what is the resulting photocurrent?
\end{enumerate}
\end{exercice}

\begin{exercice}[Wave theory vs photon model]
Explain, with reference to the photon model, why:
\begin{enumerate}
\item No electrons are emitted below the threshold frequency, however intense the light.
\item Emission of photoelectrons is instantaneous, even at very low light intensity.
\item The maximum kinetic energy of photoelectrons depends linearly on the frequency of the incident light, not on its intensity.
\end{enumerate}
\end{exercice}

\courssub{I prepare for the GCE}

\begin{exercice}[Structured calculation: Photon flux from a lamp (GCE Paper 2 style)]
A \(100\ \mathrm{W}\) sodium vapour lamp operates at an efficiency of \(15\%\) in converting electrical power to light of wavelength \(589\ \mathrm{nm}\). A photoelectric cell of cathode area \(2.0\ \mathrm{cm}^2\) is placed \(1.5\ \mathrm{m}\) from the lamp. The work function of the cathode material is \(2.14\ \mathrm{eV}\). Assume the light is emitted uniformly in all directions and that each photon incident on the cathode ejects at most one electron.
\begin{enumerate}
\item Calculate the energy of a single photon of this light in joules. [2]
\item Determine the number of photons emitted by the lamp per second. [2]
\item Find the intensity of the light at the photocell. [2]
\item Estimate the number of photons striking the cathode per second. [2]
\item Calculate the maximum possible photocurrent. [2]
\item If the actual measured photocurrent is only \(60\%\) of this maximum, suggest two reasons for the discrepancy. [2]
\end{enumerate}
\end{exercice}

\begin{exercice}[Experimental determination of Planck's constant]
A student performs a photoelectric effect experiment using a mercury vapour lamp and a set of interference filters to select different wavelengths. The stopping potential \(V_s\) is measured for each wavelength. The data obtained are:

\begin{center}
\begin{tabular}{|c|c|c|c|c|}
\hline
Wavelength \(\lambda\) (nm) & 365 & 405 & 436 & 546 \\
\hline
Stopping potential \(V_s\) (V) & 1.85 & 1.42 & 1.20 & 0.45 \\
\hline
\end{tabular}
\end{center}

\begin{enumerate}
\item Convert each wavelength to its corresponding frequency \(f\). [2]
\item Plot a graph of \(V_s\) against \(f\). [3]
\item Determine the gradient of the graph and use it to calculate a value for Planck's constant \(h\). [3]
\item Determine the work function \(\phi\) of the cathode material in electronvolts. [2]
\item The accepted value of \(h\) is \(6.63 \times 10^{-34}\ \mathrm{J\,s}\). Calculate the percentage difference between the experimental value and the accepted value. [1]
\item State two precautions the student should take to obtain accurate results. [2]
\end{enumerate}
\end{exercice}

\begin{exercice}[Essay: Wave theory vs photon model]
"The photoelectric effect cannot be explained by the wave theory of light." Discuss this statement, citing at least three experimental observations and showing how the photon model accounts for each. [15 marks]
\end{exercice}

\newpage

\coursec{Solutions to the exercises}

\begin{corrige}[Exercise 1 — MCQ: Photoelectric effect basics]
The correct answer is \textbf{(b)}.

\textbf{Justification of each option:}
\begin{itemize}
\item \textbf{(a) False.} According to Einstein's equation $E_{k,\max} = hf - \phi$, the maximum kinetic energy depends only on the \cle{frequency} $f$ and the work function $\phi$. Increasing the intensity increases the \emph{number} of emitted electrons (the photocurrent), not their individual energy.
\item \textbf{(b) True.} If $f < f_0 = \phi/h$, then each photon carries an energy $hf < \phi$, insufficient to extract even one electron. Since the photon–electron interaction is a one-to-one process, no intensity (i.e.\ no number of photons) can compensate: no electron is emitted, however bright the light.
\item \textbf{(c) False.} From $eV_s = hf - \phi$, the stopping potential depends on frequency, not on intensity.
\item \textbf{(d) False.} The photocurrent depends on the number of photons absorbed per second; for a given intensity, changing the frequency changes the photon flux, and the quantum efficiency also varies with frequency. The current is therefore not independent of $f$.
\end{itemize}
\end{corrige}

\begin{corrige}[Exercise 2 — True/False with justification]
\begin{enumerate}
\item \textbf{False.} The maximum kinetic energy is fixed by $E_{k,\max} = hf - \phi$; it depends on frequency only. Greater intensity means more photons per second, hence \emph{more} electrons emitted per second (larger photocurrent), but each with the same maximum energy.
\item \textbf{False.} The photoelectric effect demonstrates the \cle{particle (corpuscular) nature} of light: energy arrives in discrete quanta (photons). The wave theory fails to explain the threshold frequency and the instantaneous emission.
\item \textbf{False.} The threshold frequency $f_0 = \phi/h$ is a property of the \emph{metal} alone (through its work function $\phi$). It is completely independent of the intensity of the incident light.
\item \textbf{True.} In Einstein's equation $hf = \phi + E_{k,\max}$, the term $hf$ is precisely the energy of one incident photon of frequency $f$, where $h$ is Planck's constant.
\end{enumerate}
\end{corrige}

\begin{corrige}[Exercise 3 — Definitions and key concepts]
\begin{enumerate}
\item \textbf{Work function $\phi$:} the minimum energy required to remove an electron from the surface of a metal. It is characteristic of the metal and is usually expressed in electronvolts (eV).
\item \textbf{Threshold frequency $f_0$:} the minimum frequency of incident radiation below which no photoelectron is emitted, whatever the intensity. It is related to the work function by $f_0 = \phi/h$.
\item \textbf{Stopping potential $V_s$:} the minimum reverse (retarding) potential difference that must be applied between cathode and anode to reduce the photocurrent to zero, i.e.\ to stop even the most energetic photoelectrons. It satisfies $eV_s = E_{k,\max}$.
\item \textbf{Photon:} a discrete quantum (packet) of electromagnetic energy. A photon of radiation of frequency $f$ carries the energy $E = hf$ and travels at the speed of light $c$; it has zero rest mass.
\end{enumerate}
\end{corrige}

\begin{corrige}[Exercise 4 — Direct calculation: Photon energy and photoelectric emission]
\textbf{Data:} $\lambda = 400\ \mathrm{nm} = 4.00\times 10^{-7}\ \mathrm{m}$; $\phi = 2.28\ \mathrm{eV}$.

\textbf{(a) Energy of one photon.}
\[ E = hf = \frac{hc}{\lambda} = \frac{(6.63\times 10^{-34}\ \mathrm{J\,s})(3.00\times 10^{8}\ \mathrm{m\,s^{-1}})}{4.00\times 10^{-7}\ \mathrm{m}} = 4.97\times 10^{-19}\ \mathrm{J}. \]
Converting to electronvolts ($1\ \mathrm{eV} = 1.60\times 10^{-19}\ \mathrm{J}$):
\[ E = \frac{4.97\times 10^{-19}\ \mathrm{J}}{1.60\times 10^{-19}\ \mathrm{J\,eV^{-1}}} = 3.11\ \mathrm{eV}. \]

\textbf{(b) Maximum kinetic energy.} Since $E > \phi$, emission occurs:
\[ E_{k,\max} = hf - \phi = 3.11\ \mathrm{eV} - 2.28\ \mathrm{eV} = 0.83\ \mathrm{eV}. \]

\textbf{(c) Stopping potential.} From $eV_s = E_{k,\max}$, with $E_{k,\max}$ expressed in eV, the stopping potential is numerically equal to the kinetic energy in eV:
\[ V_s = 0.83\ \mathrm{V}. \]
\end{corrige}

\begin{corrige}[Exercise 5 — Threshold wavelength and work function]
\textbf{Data:} $\lambda_0 = 540\ \mathrm{nm} = 5.40\times 10^{-7}\ \mathrm{m}$.

\textbf{(a) Work function.}
\[ \phi = hf_0 = \frac{hc}{\lambda_0} = \frac{(6.63\times 10^{-34}\ \mathrm{J\,s})(3.00\times 10^{8}\ \mathrm{m\,s^{-1}})}{5.40\times 10^{-7}\ \mathrm{m}} = 3.68\times 10^{-19}\ \mathrm{J} = \frac{3.68\times 10^{-19}}{1.60\times 10^{-19}}\ \mathrm{eV} = 2.30\ \mathrm{eV}. \]

\textbf{(b) Red light at $700\ \mathrm{nm}$.} The photon energy is
\[ E = \frac{hc}{\lambda} = \frac{(6.63\times 10^{-34})(3.00\times 10^{8})}{7.00\times 10^{-7}} = 2.84\times 10^{-19}\ \mathrm{J} = 1.78\ \mathrm{eV}. \]
Since $E = 1.78\ \mathrm{eV} < \phi = 2.30\ \mathrm{eV}$ (equivalently $\lambda = 700\ \mathrm{nm} > \lambda_0 = 540\ \mathrm{nm}$), \textbf{no electron is emitted}, whatever the intensity of the red light.

\textbf{(c) Ultraviolet light at $300\ \mathrm{nm}$.}
\[ E = \frac{hc}{\lambda} = \frac{(6.63\times 10^{-34})(3.00\times 10^{8})}{3.00\times 10^{-7}} = 6.63\times 10^{-19}\ \mathrm{J} = 4.14\ \mathrm{eV}. \]
\[ E_{k,\max} = E - \phi = 4.14\ \mathrm{eV} - 2.30\ \mathrm{eV} = 1.84\ \mathrm{eV}. \]
\end{corrige}

\begin{corrige}[Exercise 6 — Einstein's equation: Graphical analysis]
\textbf{(a) Expected graph.} Einstein's equation gives
\[ eV_s = hf - \phi \quad\Longrightarrow\quad V_s = \frac{h}{e}\,f - \frac{\phi}{e}, \]
which is a straight line of gradient $h/e$ and vertical intercept $-\phi/e$. Plotting the points $(5.0,\ 0.45)$, $(7.0,\ 1.28)$, $(9.0,\ 2.10)$ (with $f$ in units of $10^{14}\ \mathrm{Hz}$) gives points that are essentially collinear.

\textbf{(b) Deduction of $h$ and $\phi$.} Using the two extreme points:
\[ \text{gradient} = \frac{\Delta V_s}{\Delta f} = \frac{2.10 - 0.45}{(9.0 - 5.0)\times 10^{14}\ \mathrm{Hz}} = \frac{1.65}{4.0\times 10^{14}} = 4.13\times 10^{-15}\ \mathrm{V\,s}. \]
Since gradient $= h/e$:
\[ h = e \times \text{gradient} = (1.60\times 10^{-19}\ \mathrm{C})(4.13\times 10^{-15}\ \mathrm{V\,s}) = 6.60\times 10^{-34}\ \mathrm{J\,s}. \]
For the work function, use one point, e.g.\ $(9.0\times 10^{14}\ \mathrm{Hz};\ 2.10\ \mathrm{V})$:
\[ \phi = hf - eV_s = (6.60\times 10^{-34}\ \mathrm{J\,s})(9.0\times 10^{14}\ \mathrm{Hz}) - (1.60\times 10^{-19}\ \mathrm{C})(2.10\ \mathrm{V}) \]
\[ \phi = 5.94\times 10^{-19}\ \mathrm{J} - 3.36\times 10^{-19}\ \mathrm{J} = 2.58\times 10^{-19}\ \mathrm{J} = 1.61\ \mathrm{eV}. \]
(Equivalently, $\phi/e$ is the magnitude of the vertical intercept, or $hf_0$ where $f_0 \approx 3.9\times 10^{14}\ \mathrm{Hz}$ is the horizontal intercept.)

\textbf{(c) Comparison.} The experimental value $6.60\times 10^{-34}\ \mathrm{J\,s}$ differs from the accepted value $6.63\times 10^{-34}\ \mathrm{J\,s}$ by less than $0.5\ \%$ — an excellent agreement. Possible sources of error: uncertainty in reading small stopping potentials, contact potentials between electrodes, impurities or oxidation of the cathode surface, impurity of the monochromatic filters, and leakage currents in the measuring circuit.
\end{corrige}

\begin{corrige}[Exercise 7 — Photon flux and photocurrent]
\textbf{Data:} $P = 60\ \mathrm{W}$; visible fraction $2\ \%$; $\lambda = 550\ \mathrm{nm}$; $A = 1\ \mathrm{cm}^2 = 1.0\times 10^{-4}\ \mathrm{m}^2$; $r = 2\ \mathrm{m}$.

\textbf{(a) Photons emitted per second.} Visible power: $P_v = 0.02 \times 60 = 1.2\ \mathrm{W}$. Energy of one photon:
\[ E_{\text{ph}} = \frac{hc}{\lambda} = \frac{(6.63\times 10^{-34}\ \mathrm{J\,s})(3.00\times 10^{8}\ \mathrm{m\,s^{-1}})}{5.50\times 10^{-7}\ \mathrm{m}} = 3.62\times 10^{-19}\ \mathrm{J}. \]
\[ N = \frac{P_v}{E_{\text{ph}}} = \frac{1.2\ \mathrm{W}}{3.62\times 10^{-19}\ \mathrm{J}} = 3.3\times 10^{18}\ \mathrm{s^{-1}}. \]

\textbf{(b) Intensity at the photocell.} The light spreads over a sphere of area $4\pi r^2$:
\[ I = \frac{P_v}{4\pi r^2} = \frac{1.2\ \mathrm{W}}{4\pi (2\ \mathrm{m})^2} = \frac{1.2}{50.3}\ \mathrm{W\,m^{-2}} = 2.4\times 10^{-2}\ \mathrm{W\,m^{-2}}. \]

\textbf{(c) Photons striking the cell per second.} Power received: $P_{\text{cell}} = IA = (2.4\times 10^{-2}\ \mathrm{W\,m^{-2}})(1.0\times 10^{-4}\ \mathrm{m}^2) = 2.4\times 10^{-6}\ \mathrm{W}$.
\[ N_{\text{cell}} = \frac{P_{\text{cell}}}{E_{\text{ph}}} = \frac{2.4\times 10^{-6}\ \mathrm{W}}{3.62\times 10^{-19}\ \mathrm{J}} = 6.6\times 10^{12}\ \mathrm{s^{-1}}. \]

\textbf{(d) Photocurrent.} With one electron per photon:
\[ I_{\text{photo}} = N_{\text{cell}}\, e = (6.6\times 10^{12}\ \mathrm{s^{-1}})(1.60\times 10^{-19}\ \mathrm{C}) = 1.1\times 10^{-6}\ \mathrm{A} \approx 1.1\ \mu\mathrm{A}. \]
\end{corrige}

\begin{corrige}[Exercise 8 — Wave theory vs photon model]
\begin{enumerate}
\item \textbf{Threshold frequency.} In the photon model, light of frequency $f$ consists of photons each of energy $hf$. An electron is extracted only if a single photon gives it at least the work function: $hf \geq \phi$. If $f < f_0 = \phi/h$, each individual photon is too ``weak''; because the interaction is one photon–one electron, photons cannot combine their energies. Hence no emission occurs, however many photons arrive per second (however intense the light). The wave theory, by contrast, predicts that a sufficiently intense wave (of any frequency) should eventually supply enough energy — which is never observed.

\item \textbf{Instantaneous emission.} A photon delivers its whole energy $hf$ to a single electron in a single collision. As soon as one photon with $hf \geq \phi$ strikes the metal, an electron can leave: there is no accumulation time, so emission is immediate even at very low intensity. The wave theory would require the electron to absorb energy gradually from the wavefront, predicting a measurable time delay at low intensity — contrary to observation.

\item \textbf{Linear dependence on frequency.} Energy conservation for one photon and one electron gives Einstein's equation $E_{k,\max} = hf - \phi$: the surplus energy of the photon above the work function becomes the electron's kinetic energy. Thus $E_{k,\max}$ grows \emph{linearly} with $f$ (gradient $h$). Increasing the intensity only increases the \emph{number} of photons per second, hence the number of emitted electrons (the current), but leaves the energy of each electron unchanged.
\end{enumerate}
\end{corrige}

\begin{corrige}[Exercise 9 — Structured calculation: Photon flux from a lamp (GCE Paper 2 style)]
\textbf{Data:} $P = 100\ \mathrm{W}$; efficiency $15\ \%$; $\lambda = 589\ \mathrm{nm}$; $A = 2.0\ \mathrm{cm}^2 = 2.0\times 10^{-4}\ \mathrm{m}^2$; $r = 1.5\ \mathrm{m}$; $\phi = 2.14\ \mathrm{eV}$.

\textbf{(a) Photon energy [2 marks].}
\[ E_{\text{ph}} = \frac{hc}{\lambda} = \frac{(6.63\times 10^{-34}\ \mathrm{J\,s})(3.00\times 10^{8}\ \mathrm{m\,s^{-1}})}{5.89\times 10^{-7}\ \mathrm{m}} = 3.38\times 10^{-19}\ \mathrm{J}. \]
\emph{Mark scheme: 1 mark for $E = hc/\lambda$; 1 mark for the value with unit.}

\textbf{(b) Photons emitted per second [2 marks].} Light power: $P_L = 0.15 \times 100 = 15\ \mathrm{W}$.
\[ N = \frac{P_L}{E_{\text{ph}}} = \frac{15\ \mathrm{W}}{3.38\times 10^{-19}\ \mathrm{J}} = 4.4\times 10^{19}\ \mathrm{s^{-1}}. \]
\emph{1 mark for the useful power $15\ \mathrm{W}$; 1 mark for $N$.}

\textbf{(c) Intensity at the cell [2 marks].}
\[ I = \frac{P_L}{4\pi r^2} = \frac{15\ \mathrm{W}}{4\pi (1.5\ \mathrm{m})^2} = \frac{15}{28.3}\ \mathrm{W\,m^{-2}} = 0.53\ \mathrm{W\,m^{-2}}. \]
\emph{1 mark for the spherical geometry $4\pi r^2$; 1 mark for the value.}

\textbf{(d) Photons striking the cathode per second [2 marks].} Power received: $P_c = IA = 0.53\ \mathrm{W\,m^{-2}} \times 2.0\times 10^{-4}\ \mathrm{m}^2 = 1.06\times 10^{-4}\ \mathrm{W}$.
\[ N_c = \frac{P_c}{E_{\text{ph}}} = \frac{1.06\times 10^{-4}\ \mathrm{W}}{3.38\times 10^{-19}\ \mathrm{J}} = 3.1\times 10^{14}\ \mathrm{s^{-1}}. \]
\emph{1 mark for $P_c = IA$; 1 mark for $N_c$.}

\textbf{(e) Maximum photocurrent [2 marks].} First check emission is possible: $E_{\text{ph}} = 3.38\times 10^{-19}\ \mathrm{J} = 2.11\ \mathrm{eV}$. This is slightly \emph{less} than $\phi = 2.14\ \mathrm{eV}$; strictly, with the rounded constants used, the photon energy is marginally below the work function. In the intended examination question (where $E_{\text{ph}} \approx 2.11\ \mathrm{eV} \approx \phi$), emission is taken as possible, and the expected answer is:
\[ I_{\max} = N_c\, e = (3.1\times 10^{14}\ \mathrm{s^{-1}})(1.60\times 10^{-19}\ \mathrm{C}) = 5.0\times 10^{-5}\ \mathrm{A} = 50\ \mu\mathrm{A}. \]
\emph{1 mark for $I = N_c e$; 1 mark for the value. (A candidate who remarks that $E_{\text{ph}}$ is marginally below $\phi$ and concludes that the current is essentially zero demonstrates excellent understanding and should receive full credit.)}

\textbf{(f) Reasons for the measured current being only $60\ \%$ of the maximum [2 marks].} Any two of:
\begin{itemize}
\item Quantum efficiency below $100\ \%$: not every absorbed photon ejects an electron (some photons are reflected, some electrons lose their energy in collisions inside the metal before escaping).
\item Some emitted electrons are not collected by the anode (recombination, unfavourable geometry, insufficient collecting voltage).
\item Absorption/scattering of light by the glass envelope and the air; cathode surface contamination or oxidation increasing the effective work function.
\end{itemize}
\emph{1 mark per valid reason, maximum 2.}
\end{corrige}

\begin{corrige}[Exercise 10 — Experimental determination of Planck's constant]
\textbf{(a) Frequencies [2 marks].} Using $f = c/\lambda$:

\begin{center}
\begin{tabular}{|c|c|c|c|c|}
\hline
\rowcolor{popblueL} $\lambda$ (nm) & 365 & 405 & 436 & 546 \\
\hline
$f = c/\lambda$ ($\times 10^{14}$ Hz) & 8.22 & 7.41 & 6.88 & 5.49 \\
\hline
$V_s$ (V) & 1.85 & 1.42 & 1.20 & 0.45 \\
\hline
\end{tabular}
\end{center}
\emph{1 mark for the formula, 1 mark for the four correct values.}

\textbf{(b) Graph [3 marks].} Plot $V_s$ (vertical axis) against $f$ (horizontal axis). The four points are nearly collinear; draw the best-fit straight line. \emph{Marks: correct axes with labels and units (1); correct plotting (1); best-fit line (1).}

\textbf{(c) Gradient and Planck's constant [3 marks].} Using the extreme points:
\[ \text{gradient} = \frac{1.85 - 0.45}{(8.22 - 5.49)\times 10^{14}\ \mathrm{Hz}} = \frac{1.40}{2.73\times 10^{14}} = 5.13\times 10^{-15}\ \mathrm{V\,s}. \]
Since $V_s = \dfrac{h}{e}f - \dfrac{\phi}{e}$, the gradient equals $h/e$:
\[ h = e \times \text{gradient} = (1.60\times 10^{-19}\ \mathrm{C})(5.13\times 10^{-15}\ \mathrm{V\,s}) = 8.2\times 10^{-34}\ \mathrm{J\,s}. \]
\emph{Marks: gradient from a large triangle on the line (1); relation gradient $= h/e$ (1); value of $h$ with unit (1).}

\textbf{(d) Work function [2 marks].} Using the point $(8.22\times 10^{14}\ \mathrm{Hz};\ 1.85\ \mathrm{V})$:
\[ \phi = hf - eV_s = (8.2\times 10^{-34}\ \mathrm{J\,s})(8.22\times 10^{14}\ \mathrm{Hz}) - (1.60\times 10^{-19}\ \mathrm{C})(1.85\ \mathrm{V}) \]
\[ \phi = 6.74\times 10^{-19}\ \mathrm{J} - 2.96\times 10^{-19}\ \mathrm{J} = 3.78\times 10^{-19}\ \mathrm{J} \approx 2.4\ \mathrm{eV}. \]
\emph{1 mark for the method (intercept or point substitution), 1 mark for the value in eV.}

\textbf{(e) Percentage difference [1 mark].}
\[ \frac{|8.2 - 6.63|}{6.63}\times 100\ \% \approx 24\ \%. \]
The experimental value is about $24\ \%$ above the accepted value — a large but plausible discrepancy for a school experiment with only four points and simple apparatus.

\textbf{(f) Precautions [2 marks].} Any two of:
\begin{itemize}
\item Use filters of narrow bandwidth to ensure truly monochromatic light, and shield the cell from stray light.
\item Clean (or avoid oxidation of) the cathode surface so that the work function is well defined.
\item Measure the stopping potential accurately: adjust the retarding voltage slowly until the current just falls to zero, using a sensitive galvanometer/nanoammeter.
\item Account for (or minimise) contact potential differences between the electrodes; repeat readings and take averages.
\end{itemize}
\emph{1 mark per valid precaution, maximum 2.}
\end{corrige}

\begin{corrige}[Exercise 11 — Essay: Wave theory vs photon model]
\textbf{Indicative mark scheme (15 marks):} statement of the conflict and of the photon hypothesis (3); at least three experimental observations with the failure of wave theory for each (3 $\times$ 2 = 6); explanation of each observation by the photon model (3 $\times$ 2 = 6).

\textbf{Model answer.}

\emph{Introduction.} The classical wave theory regards light as a continuous electromagnetic wave whose energy is spread over the whole wavefront and proportional to the intensity. Applied to the photoelectric effect, it makes three clear predictions, all contradicted by experiment. Einstein's photon model (1905), in which light of frequency $f$ consists of quanta of energy $E = hf$ absorbed individually by electrons, accounts naturally for every observation.

\emph{Observation 1 — Existence of a threshold frequency.} Experiment shows that for each metal there is a frequency $f_0$ below which no electron is ever emitted, whatever the intensity or duration of illumination. The wave theory predicts that a sufficiently intense wave of \emph{any} frequency should eventually eject electrons, since energy could accumulate continuously. In the photon model, extraction requires a single photon to supply at least the work function: $hf \geq \phi$. If $f < f_0 = \phi/h$, every photon is individually too weak, and photons cannot pool their energy because absorption is a one-photon–one-electron event. Intensity only changes the number of photons, not their individual energy.

\emph{Observation 2 — Instantaneous emission.} Photoelectrons appear within less than about $10^{-9}\ \mathrm{s}$ of illumination, even at extremely low intensity. The wave theory predicts a delay: at low intensity the energy flux is small, and an electron would need a long time to gather the energy $\phi$ from the spread-out wavefront. In the photon model, a photon delivers its entire energy $hf$ in a single collision; the very first photon absorbed with $hf \geq \phi$ releases an electron immediately.

\emph{Observation 3 — Kinetic energy depends on frequency, not intensity.} Experiment (Millikan) shows that the maximum kinetic energy, measured through the stopping potential via $eV_s = E_{k,\max}$, increases \emph{linearly} with frequency and is completely unaffected by intensity; intensity only changes the photocurrent. The wave theory links energy to intensity, so brighter light should give faster electrons — it does not. The photon model gives Einstein's equation
\[ hf = \phi + E_{k,\max}, \]
a linear relation in $f$ with gradient $h$, exactly as observed; raising the intensity merely raises the number of photons, hence the number of electrons per second.

\emph{(Additional supporting observation — photocurrent proportional to intensity.} For $f > f_0$, the saturation current is proportional to the intensity, which the photon model interprets directly: more photons per second $\Rightarrow$ more electrons per second.\emph{)}

\emph{Conclusion.} Because the wave theory fails on the threshold frequency, the absence of any time delay, and the dependence of electron energy on frequency rather than intensity, the photoelectric effect cannot be explained by it. The photon model explains all observations quantitatively through $E_{k,\max} = hf - \phi$, establishing that light exchanges energy with matter in discrete quanta — a cornerstone of quantum physics and of wave–particle duality.
\end{corrige}

\newpage

\coursec{Summary sheet}

\begin{popfigure}
\begin{tikzpicture}[grow cyclic, mindmap, concept color=popblueL,
  level 1 concept/.append style={sibling angle=60, level distance=4.1cm, minimum size=2.5cm, font=\small},
  level 2 concept/.append style={sibling angle=45, level distance=2.5cm, minimum size=1.9cm, font=\scriptsize},
  every node/.style={concept, circular drop shadow}]
\node[concept, concept color=popblue, text=white, font=\bfseries\small, minimum size=2.9cm]{Photoelectric Effect}
  child[concept color=poporangeL]{ node[concept]{Observation} 
      child{ node[concept]{Emission of $e^-$ from metal} }
      child{ node[concept]{Instantaneous, threshold $\nu_0$} }
  }
  child[concept color=popgreenL]{ node[concept]{Failure of wave theory}
      child{ node[concept]{$I$ should control $E_{c\max}$} }
      child{ node[concept]{No delay predicted} }
  }
  child[concept color=poppinkL]{ node[concept]{Photons}
      child{ node[concept]{$E=h\nu=\dfrac{hc}{\lambda}$} }
      child{ node[concept]{$h=\SI{6.63e-34}{J.s}$} }
  }
  child[concept color=popgold!60]{ node[concept]{Einstein equation}
      child{ node[concept]{$h\nu=W_0+E_{c\max}$} }
      child{ node[concept]{$W_0=h\nu_0$} }
      child{ node[concept]{$E_{c\max}=eU_0$} }
  }
  child[concept color=popgreen!50]{ node[concept]{Duality}
      child{ node[concept]{$\lambda=\dfrac{h}{p}$ de Broglie} }
      child{ node[concept]{Interference vs impact} }
  }
  child[concept color=popblueL]{ node[concept]{Applications}
      child{ node[concept]{Photocells, solar panels} }
      child{ node[concept]{Electron microscope} }
  };
\end{tikzpicture}
\legende{Mind map of the chapter: the photoelectric effect, photons, Einstein's equation and wave--particle duality.}
\end{popfigure}

\begin{bilan}
\textbf{Key definitions}
\begin{itemize}
  \item \cle{Photoelectric effect}: emission of electrons from a metal surface when it is illuminated by electromagnetic radiation of sufficiently high frequency.
  \item \cle{Photon}: quantum (packet) of light energy; each photon carries $E=h\nu$.
  \item \cle{Planck constant}: $h=\SI{6.63e-34}{J.s}$, linking energy and frequency.
  \item \cle{Work function} $W_0$: minimum energy needed to extract an electron from the metal; $W_0=h\nu_0=\dfrac{hc}{\lambda_0}$.
  \item \cle{Threshold frequency} $\nu_0$ (threshold wavelength $\lambda_0$): below $\nu_0$ (above $\lambda_0$) no electron is emitted, whatever the intensity.
  \item \cle{Stopping potential} $U_0$: reverse voltage that stops the most energetic photoelectrons; $E_{c\max}=eU_0$.
\end{itemize}

\textbf{Laws and formulas to know}
\begin{itemize}
  \item Photon energy: $E=h\nu=\dfrac{hc}{\lambda}$ with $c=\SI{3.00e8}{m/s}$.
  \item Useful conversion: $\SI{1}{eV}=\SI{1.60e-19}{J}$; $E(\mathrm{eV})\approx \dfrac{1240}{\lambda(\mathrm{nm})}$.
  \item \cle{Einstein's photoelectric equation}: $h\nu = W_0 + E_{c\max}$, i.e. $E_{c\max}=h\nu-h\nu_0$.
  \item Threshold: $\nu_0=\dfrac{W_0}{h}$, $\lambda_0=\dfrac{hc}{W_0}$.
  \item Stopping potential: $eU_0=E_{c\max}=h(\nu-\nu_0)$; the graph of $U_0$ against $\nu$ is a straight line of slope $h/e$.
  \item Saturation current is proportional to light intensity (number of photons per second).
  \item de Broglie wavelength: $\lambda=\dfrac{h}{p}=\dfrac{h}{mv}$; for an electron accelerated through $U$: $\lambda=\dfrac{h}{\sqrt{2meU}}$.
\end{itemize}

\textbf{Methods}
\begin{itemize}
  \item Always convert $\lambda$ to metres and energies to joules before calculating; convert back to eV at the end if asked.
  \item To test whether emission occurs: compare $\nu$ with $\nu_0$ (or $\lambda$ with $\lambda_0$, or $h\nu$ with $W_0$).
  \item To find $W_0$ or $h$ experimentally: plot $E_{c\max}$ (or $U_0$) versus $\nu$; slope gives $h$ (or $h/e$), intercept gives $-W_0$.
  \item Energy balance reasoning: photon energy = extraction energy + kinetic energy of the fastest electron.
\end{itemize}

\textbf{Common mistakes}
\begin{itemize}
  \item Confusing intensity with frequency: intensity changes the \emph{number} of electrons, never their maximum kinetic energy.
  \item Mixing up $\nu$ and $\lambda$ in the threshold condition: $\nu>\nu_0$ but $\lambda<\lambda_0$.
  \item Forgetting to convert nm to m or eV to J.
  \item Writing $E_{c\max}=h\nu$ and forgetting $W_0$.
  \item Believing a more intense red light can eject electrons when $\nu<\nu_0$: it cannot.
  \item Using $\lambda=h/(mv)$ with $v$ in km\,s$^{-1}$ or $m$ in grams.
\end{itemize}

\textbf{GCE tips}
\begin{itemize}
  \item Quote Einstein's equation and define every symbol before substituting values.
  \item State clearly why wave theory fails: emission is instantaneous and depends on frequency, not intensity.
  \item In graph questions, read the slope carefully with units; $h/e$ has units V\,s.
  \item Give numerical answers with SI units and 3 significant figures.
  \item For duality essays: light shows particle behaviour in the photoelectric effect and wave behaviour in interference/diffraction; matter shows wave behaviour through electron diffraction (de Broglie).
\end{itemize}
\end{bilan}

\findecours

\end{document}
